% Modernised reading — fonds Alexandre Grothendieck, shelfmark 145, pages 1-37.
%
% This is an INTERPRETATION, not a transcription. It re-reads the pages in
% today's notation and vocabulary, states the hypotheses the manuscript leaves
% implicit, and drops the critical apparatus so that the mathematics reads
% straight through. Everything the brackets and underlines carried moves into a
% footnote. It is held to being mathematically correct as it stands; where the
% manuscript is loose or mistaken the true statement appears and the footnote
% says what the page has. In any dispute of substance the transcription
% governs: whoever cites Grothendieck cites batch-01.fr.tex and batch-02.fr.tex.
%
% Scope: the folder is thirty-seven pages in two batches (1-20, 21-37), both
% transcribed, read whole before any of this was written, and written as ONE
% file for the whole shelfmark.
%
% The spine. Three loose sheets (1-3), then two MOUTURES of one text, kept
% distinct because their difference is content:
%   - first draft, pages 4-14 (his pp. 1-3 and following): outer « loop »
%     automorphisms of the profinite pi^_{0,3}, computed in Aut(pi^) with the
%     relation l_0 l_1 l_inf = 1; ends on the quasi-square roots eps_i and a
%     struck, false relation eps_0 eps_1 eps_inf = 1;
%   - second draft, pages 14-37 (« Critique des notations et de l'approche »,
%     his pp. 7-17): the discrete extension E_{0,3} of Gamma_{0,3} = S_3 x Z/2
%     by pi_{0,3}, built geometrically on the Riemann sphere with the relation
%     l_inf l_1 l_0 = 1; then the profinite E^, the group E^^ = Aut_lac(pi^),
%     and the extension S E^^ by Z^^{0,1,inf} of the last pages.
% Page 20 (no number of his) carries (13)-(16), which follow (11)-(12) of page
% 21 (his p. 9): it is read after page 21.
%
% Notation fixed for the whole folder. pi = pi_{0,3} is the discrete free
% group, pi^ its profinite completion (the first draft's Pi^ and the Pi of pages
% 15-20 are both written pi / pi^). int(g)(x) = g x g^-1; automorphisms compose
% as maps. Autext is written Out. Aut_lac, Out_lac: automorphisms permuting the
% three inertia classes with one multiplier p. The centraliser of S_3 in
% Out_lac(pi^) is written bold Gamma throughout: page 10 calls it Pi, page 23
% bold Gamma; they are the same group. The superscript ! means « pure » (sigma =
% 1). In S E^^ the convention of page 32, u(l_i) = int(g_i)(l_{sigma(i)}^p), is
% used everywhere; pages 23-24 write int(g_i^-1) (the -1 added later) and are
% translated.
%
% Every identity in Aut(pi) asserted below as true was checked by machine in
% the free group of rank 2 (script outside the repository), in each draft's own
% convention; these checks decided the following departures.
%
% Substantive departures from the pages, each footnoted where it occurs:
%   - p. 5: the fixed points of rho(z) = 1/(1-z) are -j and -j^2 = e^{±i pi/3},
%     not « ± j »;
%   - pp. 7, 9: alpha, q, r range over Z^, not Z^*;
%   - p. 9 margin: the page has beta = r, gamma = -q; in fact beta = -r = q and
%     beta = gamma, as page 11 writes (sigma_inf(g^-1) = l_0^beta g l_1^beta),
%     so the invariant (q, r) of page 9 is one number, r = -q;
%   - pp. 13-14 (struck): eps_0 eps_1 eps_inf = 1 and sigma_inf sigma_1 sigma_0
%     = 1 are false; eps_0 eps_1 eps_inf = sigma_1 is true, which is what the
%     unstruck « Sans doute » line at the foot of page 14 reads, and its three
%     circular conjugation relations are true;
%   - p. 17 margin: sigma_0(l_0) = l_inf^-1 l_1^-1 (page 19), not l_0^-1 l_1^-1;
%   - p. 19: (4)-(6) is not a presentation of E^+ (its abelianisation is too
%     big); with rho sigma_i rho^-1 = sigma_rho(i) added it is, and E^+ =
%     <sigma_0, rho | sigma_0^2 = rho^3 = 1> = PSL_2(Z) (identification ours);
%   - p. 20: phi(S_3) is Gamma_P (the isotropy group of P), not Gamma;
%   - p. 21, (10): the fourth line must read eps_1^2 where eps_inf^2 is read;
%   - p. 25: the covariance is gamma_rho(i)(rho~ u rho~^-1) = rho(gamma_i(u)),
%     as (15) of page 36 gives, not gamma_i(...) = gamma_{i+1}(u);
%   - p. 29: sigma_0 sigma_1 = rho tau_inf tau_1 = tau_0 tau_inf rho (relation (2)
%     of page 17), not tau_inf tau_1 rho;
%   - p. 30: the relation [rho, sigma_0][rho^-1, tau_0] = 1 is false (already
%     modulo pi); the true defining relation is (sigma_0 tau_0 rho)^2 = 1, and E
%     is then the Coxeter group (2,3,inf), i.e. PGL_2(Z) (identification ours);
%   - p. 32 margin: the kernel of S E^^ -> Aut_lac(pi^) is i(Z^^3), not i(pi^);
%   - p. 35: (13) must read l_sigma(inf)^p l_sigma(1)^p l_sigma(0)^p = 1 (the
%     order of (2) and (14)); the page reverses it;
%   - p. 35, « Exercice ! »: proved here by a dihedral quotient;
%   - p. 36, (16): the automorphism applied to the g_i is sigma~_0 (the one
%     underlying sigma~'_0 of (12)), written sigma_0 on the page;
%   - p. 1: H is the commensurator of pi; « pi', pi'' of the same index » needs
%     u to preserve a finite covolume; « H_n is a subgroup » is not evident and
%     is not asserted;
%   - p. 2: the boxed relations (1), (2) are reported, not checked (the sheet
%     does not fix its conventions); the algebra that follows them is checked.
%
% Announced and never established in the folder: whether Out_lac(pi^) = S_3 x
% bold Gamma (page 23, « j'ignore »); a criterion for (p, g) to give an
% automorphism (page 8, « Ça semble délicat »); generators and relations for
% S Gamma and S E (pages 31, 36, « je ne vais pas m'y attarder »); the
% commutation modulo inner automorphisms announced on the last line of page 37.
%
% Pass: Opus 5.5 (claude-opus-5-5), 2026-09-24 — first pass, unchecked against the pages by a human.
\documentclass[11pt,a4paper]{article}
% Le préambule commun à toutes les transcriptions du fonds Grothendieck.
%
% Il tient en une idée : **l'incertitude se compose, elle ne se résout pas.**
% Une transcription qui lisse un mot douteux a détruit la seule chose pour
% laquelle on la faisait — un lecteur doit pouvoir savoir, à chaque endroit,
% ce qui était sur le feuillet et ce qui a été ajouté. Les six macros qui
% suivent sont donc l'appareil critique, et non de la décoration.
%
% Les mêmes commandes sont reconnues par `scripts/render.mjs`, qui produit la
% vue de lecture du site. Une macro ajoutée ici doit l'être aussi là-bas,
% faute de quoi le rendu échouera bruyamment — ce qui est voulu : un rendu
% silencieusement faux est pire qu'un rendu absent.

\ProvidesPackage{grothendieck}[2026/08/08 Transcriptions du fonds Grothendieck]

\usepackage{amsmath,amssymb,amsthm}
\usepackage{mathrsfs}
\usepackage{stmaryrd}
% Les diagrammes commutatifs sont partout dans ce fonds : c'est la langue
% même de la géométrie algébrique de Grothendieck, et une transcription qui
% les rendrait en prose ne transcrirait rien.
\usepackage{tikz-cd}
% Français par défaut — c'est la langue des feuillets — mais l'anglais est
% chargé aussi : la traduction et le résumé sont des documents anglais, et la
% typographie française (espaces avant les deux-points) y serait une faute.
% Ces documents basculent par \selectlanguage{english} en tête de corps.
\usepackage[english,french]{babel}
\usepackage{fontspec}
\usepackage{geometry}
\geometry{a4paper,margin=25mm,marginparwidth=28mm}
\usepackage{marginnote}
\usepackage{xcolor}
% ulem, not soul. `soul`'s \st and \ul are built on a character-by-character
% scan that XeLaTeX's font machinery breaks: the document fails with an
% undefined control sequence rather than degrading. ulem does the same two jobs
% and is Unicode-engine safe. `normalem` stops it hijacking \emph, which the
% transcriptions use for Grothendieck's own emphasis.
\usepackage[normalem]{ulem}
\usepackage{hyperref}
\hypersetup{colorlinks=true,linkcolor=black,urlcolor=black,pdfborder={0 0 0}}

% --- Métadonnées du lot -----------------------------------------------------
% Reprises telles quelles par le script de rendu, qui les affiche en tête de
% la vue de lecture. Elles fixent ce que le fichier prétend couvrir : sans
% elles, une transcription flotte sans point d'attache dans un fonds de seize
% mille feuillets.

\newcommand{\folder}[1]{\gdef\grothFolder{#1}}
\newcommand{\batch}[1]{\gdef\grothBatch{#1}}
\newcommand{\leaves}[2]{\gdef\grothFirst{#1}\gdef\grothLast{#2}}
\newcommand{\foldertitle}[1]{\gdef\grothFoldertitle{#1}}
\gdef\grothFolder{?}\gdef\grothBatch{?}\gdef\grothFirst{?}\gdef\grothLast{?}\gdef\grothFoldertitle{}

% --- Le résumé --------------------------------------------------------------
%
% Il ouvre la lecture modernisée, sous le titre « Résumé », et il est composé
% en petit. Ce n'est pas un ornement : c'est le seul passage du document qui
% ne soit pas des mathématiques, et un lecteur doit pouvoir le reconnaître
% avant de l'avoir lu — puis le sauter s'il connaît déjà le sujet, ou s'y
% arrêter s'il ne le connaît pas. Le filet qui le suit marque la reprise du
% corps.

\newenvironment{resume}
  {\small\setlength{\parskip}{0.45em}}
  {\par\medskip\hrule\medskip}

% --- Mots-clés ---------------------------------------------------------------
%
% En anglais, en clôture du résumé de la lecture modernisée : le vocabulaire
% moderne sous lequel chercher ce que les feuillets construisent. Le manifeste
% du site (`npm run manifest`) les extrait pour étiqueter la cote entière —
% cette ligne est la source unique des tags, il n'y a pas de fichier à côté.
\newcommand{\keywords}[1]{\par\smallskip\noindent
  {\footnotesize\textsc{Keywords} --- \textit{#1}}\par}

% --- L'appareil critique ----------------------------------------------------

% Le numéro de feuillet, en marge. C'est l'ancre qui permet de régler un
% désaccord sur une formule en regardant *un* feuillet plutôt qu'une chemise
% de six cents. Le site s'en sert aussi pour tourner les pages du fac-similé
% quand on fait défiler la transcription.
\newcommand{\leaf}[1]{%
  \marginnote{\footnotesize\color{gray!70}\textsf{#1}}%
  \label{leaf:#1}\ignorespaces}

% Illisible : on ne devine pas. Un mot inventé qui se lit comme les autres est
% le pire résultat possible.
\newcommand{\ill}{\textcolor{red!60!black}{[\,\ldots\,]}}

% Lecture douteuse : le mot est donné, et signalé comme incertain.
\newcommand{\uncertain}[1]{\uline{#1}}

% Ajout de l'éditeur — une lettre manquante, un mot sous-entendu.
\newcommand{\add}[1]{\textcolor{blue!50!black}{[#1]}}

% Biffé par Grothendieck lui-même. Ce qu'il a rayé fait partie du document :
% une rature dit souvent où la pensée a changé de direction.
\newcommand{\struck}[1]{\sout{#1}}

% Note du transcripteur, sur l'état matériel du feuillet.
\newcommand{\note}[1]{\par\smallskip{\small\color{gray!60!black}\textit{[#1]}}\par\smallskip}

% Note marginale de Grothendieck, distincte du corps du texte.
\newcommand{\marginal}[1]{\par\smallskip\noindent\hspace{1em}%
  {\small\color{gray!60!black}#1}\par\smallskip}

% --- Titre ------------------------------------------------------------------

\AtBeginDocument{%
  \begin{center}
    {\large\bfseries Fonds Alexandre Grothendieck --- cote n\textsuperscript{o}~\grothFolder}\\[2pt]
    {\itshape\grothFoldertitle}\\[4pt]
    {\small feuillets \grothFirst--\grothLast\ (lot \grothBatch)}\\[2pt]
    {\footnotesize\color{gray!60!black}Transcription établie d'après le fac-similé numérisé
     par l'Université de Montpellier.\\
     Les lectures incertaines sont soulignées, les illisibles notées [\,\ldots\,],
     les ajouts entre crochets.}
  \end{center}
  \vspace{1em}}

\endinput


\folder{145}
\pages{1}{37}
\dating{1981}
\watermark{Édition de démonstration\\interprétation personnelle de l'œuvre}
\foldertitle{Printemps 81 (vieille rédaction) : notes manuscrites (s.d.) — lecture modernisée du dossier entier}

\begin{document}

\section*{Résumé}

\begin{resume}
Ces trente-sept pages étudient un seul objet, le plus petit qui ne soit pas
trivial : la sphère de Riemann privée de trois points, $0$, $1$ et $\infty$.
Une fourmi qui s'y promène peut tourner autour de chacun des trous, et ses
trajets de retour au point de départ forment un groupe, le \emph{groupe
fondamental}. Ici il est libre à deux générateurs : un lacet autour de $0$, un
lacet autour de $1$ ; le lacet autour de $\infty$ s'en déduit, car faire le
tour des trois trous l'un après l'autre revient à ne rien faire.

La question du dossier est celle des symétries. La sphère trouée en a de
visibles : six transformations qui permutent les trois points (les
« permutations » $z \mapsto 1-z$, $z \mapsto 1/z$, etc.), et la conjugaison
complexe, qui la retourne comme un miroir. Chacune agit sur les lacets. Mais le
groupe des lacets, surtout quand on le « complète » en y ajoutant toutes les
limites possibles de lacets — c'est le groupe \emph{profini}, qui voit tous les
revêtements finis à la fois —, a beaucoup plus de symétries que la sphère n'en
montre. On demande seulement à ces symétries cachées de respecter la trace des
trous : envoyer un petit lacet autour d'un point sur un petit lacet autour d'un
point, éventuellement parcouru un nombre « $p$ » de fois, le même pour les
trois, que Grothendieck appelle le \emph{multiplicateur}. Qui sont-elles, et
comment s'accordent-elles avec les six permutations visibles ? C'est le
problème des pages 4 à 14 : il y écrit une symétrie cachée comme un couple
$(p, g)$, un multiplicateur et un élément du groupe, et traduit « commuter aux
permutations » en deux équations sur $g$.

Puis il s'arrête. Un calcul vient de lui donner une relation qu'il juge
« inquiétante », et elle est en effet fausse ; la page suivante s'ouvre sur une
« Critique des notations et de l'approche ». Ce qu'elle demande tient en une
phrase : ne plus calculer à l'aveugle dans le groupe complété, mais d'abord
dans le groupe ordinaire, avec la géométrie sous les yeux — la sphère, ses axes
de rotation, ses grands cercles, et des chemins tracés dessus. La seconde
rédaction (pages 14 à 37) le fait, et c'est son intérêt : chaque signe qui, dans
la première, se décidait par un calcul qu'on pouvait rater, se lit désormais
sur un dessin. Les symétries visibles et les lacets s'y assemblent en un seul
groupe discret, qu'on sait aujourd'hui être le \emph{groupe modulaire étendu}
$PGL_2(\mathbb{Z})$ ; ses sous-groupes finis correspondent aux points de la
sphère fixés par une symétrie, et le dossier les passe en revue un à un. On y
croise, sous le nom de « groupe cartographique de Malgoire », le groupe qui
décrit les cartes tracées sur une surface. Les dernières pages reviennent au
groupe complété, avec un dispositif plus fin : chaque symétrie cachée y est
accompagnée d'un choix de « rectificateurs », un par trou, et l'on obtient une
description en coordonnées où il ne reste qu'une équation à écrire. Le dossier
s'achève sur un exercice qu'il se donne, et que nous résolvons.

Le dossier ne dit pas pourquoi ces symétries l'intéressent, mais le titre de la
chemise — « Printemps 81 » — le situe : c'est le moment de \emph{La Longue
Marche à travers la théorie de Galois}. Le groupe de Galois des nombres
algébriques agit sur ce groupe de lacets complété, et cette action est
fidèle : c'est le théorème de Belyi (1979). Chaque symétrie arithmétique est
donc une des symétries cachées qu'on cherche ici, avec pour multiplicateur son
action sur les racines de l'unité. Mesurer le groupe de toutes ces symétries
cachées, c'est mesurer, de l'extérieur, le groupe de Galois lui-même. Dix ans
plus tard, Drinfeld et Ihara en ont fait le \emph{groupe de
Grothendieck--Teichmüller} ; les équations des pages 6 à 11 sont de la famille
des deux premières de ses trois relations de définition.

Pour aller plus loin : groupe fondamental de la droite projective moins trois
points, automorphismes préservant l'inertie, groupe de
Grothendieck--Teichmüller, groupe modulaire étendu, groupe fondamental
orbifold, groupe cartographique, commensurateur.

\keywords{thrice-punctured sphere, free profinite group, inertia-preserving automorphism, outer automorphism group, Grothendieck–Teichmüller group, anharmonic group, orbifold fundamental group, extended modular group, cartographic group, group extension, crossed homomorphism, commensurator}
\end{resume}

\pagerange{1}{37}
\section*{Le fil du dossier, et les conventions}

\subsection*{Les stations}

\begin{itemize}
\item \textbf{Trois feuillets épars (pages 1 à 3)}, sur trois papiers
différents, sans pagination de sa main : le groupe $H$ des éléments qui
descendent en isomorphismes entre revêtements finis (un commensurateur,
page 1) ; les automorphismes à lacets commutant à $\sigma_0$ (page 2) ; une
page barrée sur les espaces classifiants et les automorphismes compatibles à
un morphisme (page 3).
\item \textbf{Première rédaction (pages 4 à 14, ses pages 1 à 3 et suivantes)} :
« Automorphismes extérieurs des groupes à lacets profinis ». L'action
extérieure de $\mathfrak{S}_3 \times \mathbb{Z}/2$ sur $\widehat{\pi}$ ; les
automorphismes extérieurs commutant à $\rho$, puis à $\mathfrak{S}_3$, décrits
par $(p, g)$ ; les racines quasi-carrées $\varepsilon_i$.
\item \textbf{Seconde rédaction, le groupe discret (pages 14 à 22)} : la
critique de la page 14 ; la sphère et ses symétries ; l'extension
$\mathcal{E} = \mathcal{E}_{0,3}$ « à la Malgoire » ; sa présentation ; le
relèvement canonique du stabilisateur de $P$ ; le groupe cartographique.
\item \textbf{Le complété et les automorphismes à lacets (pages 22 à 25)} :
$\widehat{\mathcal{E}} \subset \hat{\hat{\mathcal{E}}} =
\mathrm{Aut}_{\mathrm{lac}}(\widehat{\pi})$, le centralisateur
$\boldsymbol{\Gamma}$ de $\mathfrak{S}_3$, la première forme de
$S\hat{\hat{\mathcal{E}}}$.
\item \textbf{Scindages partiels et présentations de $\mathcal{E}$ (pages 26 à
31)} : les sous-groupes finis, points fixes par points fixes ; une
présentation par trois générateurs.
\item \textbf{Le groupe $S\hat{\hat{\mathcal{E}}}$ (pages 32 à 37, ses pages 15 à
17)} : description en coordonnées, loi de groupe, le sous-groupe des
« rectificateurs » triviaux, l'exercice de la page 35, les centralisateurs de
$\tilde\rho$ et $\tilde\sigma_0$.
\end{itemize}

Les deux rédactions ne sont pas deux versions à fusionner. La première
calcule dans $\mathrm{Aut}(\widehat{\pi})$ avec la relation
$l_0 l_1 l_\infty = 1$ ; la seconde construit un groupe discret à partir de
chemins sur la sphère, et la même géométrie lui impose l'ordre opposé,
$l_\infty l_1 l_0 = 1$ — il le souligne (« et non $l_0 l_1 l_\infty = 1$ ! »).
Les éléments $\sigma_i$, $\rho$, $\varepsilon_i$ ne sont donc pas les mêmes
d'une rédaction à l'autre, et chacune est lue ici dans ses propres
conventions.

La page 20, sans numéro de sa main, porte les formules (13) à (16) ; elles
suivent (11) et (12), qui sont à la page 21 (sa page 9). On lit donc la page 20
après la page 21.

\subsection*{Les conventions}

\emph{Groupes.} $\pi = \pi_{0,3} = \pi_1(\mathbb{P}^1_{\mathbb{C}} \setminus
\{0,1,\infty\})$ est le groupe discret, libre sur deux des trois lacets
$l_0, l_1, l_\infty$ ; $\widehat{\pi}$ est son complété profini, libre profini
de rang $2$\footnote{Le dossier écrit $\hat\Pi$ et $\hat\pi$ indifféremment aux
pages 2 à 14 pour le groupe profini, et $\Pi$ pour le groupe discret aux pages
15 à 20 ; à partir de la page 21, $\pi$ et $\hat\pi$. On s'en tient à $\pi$ et
$\widehat{\pi}$.}. $L_i$ désigne le sous-groupe (fermé) engendré par $l_i$,
isomorphe à $\widehat{\mathbb{Z}}$ dans $\widehat{\pi}$. On écrit
$\mathrm{int}(g)(x) = g x g^{-1}$, et les automorphismes se composent comme des
applications ; comme $\widehat{\pi}$ est de centre trivial, $g \mapsto
\mathrm{int}(g)$ identifie $\widehat{\pi}$ à un sous-groupe de
$\mathrm{Aut}(\widehat{\pi})$, et l'on écrit souvent $g$ pour $\mathrm{int}(g)$,
comme le fait le dossier.

\emph{Symétries visibles.} $\mathfrak{S}_3$ est le groupe des permutations de
$\{0, 1, \infty\}$ ; $\rho$ est la permutation circulaire $0 \mapsto 1 \mapsto
\infty \mapsto 0$, et $\sigma_i$ la transposition qui fixe $i$. Sur la sphère,
$\rho(z) = 1/(1-z)$, $\sigma_0(z) = z/(z-1)$, $\sigma_1(z) = 1/z$,
$\sigma_\infty(z) = 1-z$, et $\tau(z) = \bar z$. On pose $\Gamma =
\Gamma_{0,3} = \mathfrak{S}_3 \times \mathbb{Z}/2$, le second facteur engendré
par $\tau$ ; quand il faut distinguer un élément de $\Gamma$ de ses relèvements,
on le pointe : $\dot\sigma_i$, $\dot\rho$.

\emph{Automorphismes à lacets.} Un automorphisme $u$ de $\widehat{\pi}$ est
\emph{à lacets} s'il existe $\sigma \in \mathfrak{S}_3$ et $p \in
\widehat{\mathbb{Z}}^{*}$ tels que $u(l_i)$ soit conjugué à $l_{\sigma(i)}^{p}$
pour $i = 0, 1, \infty$ ; $p = \chi(u)$ est son \emph{multiplicateur}. On note
$\mathrm{Aut}_{\mathrm{lac}}(\widehat{\pi})$ leur groupe, et
$\mathrm{Out}_{\mathrm{lac}}(\widehat{\pi})$ son image dans
$\mathrm{Out}(\widehat{\pi})$\footnote{Le dossier écrit « Aut ext » (ici
$\mathrm{Out}$), avec l'indice « lac », et ne définit nulle part d'un seul
tenant les automorphismes « à lacets » ; la définition ci-dessus est celle que
les formules (1) à (3) de la page 32 rendent explicite. Qu'un même exposant $p$
serve pour les trois lacets n'est pas une hypothèse : si $u(l_i)$ est conjugué
à $l_{\sigma(i)}^{p_i}$, l'abélianisé $\widehat{\mathbb{Z}}^2$, où $\bar l_0 +
\bar l_1 + \bar l_\infty = 0$ et où $\bar l_0, \bar l_1$ forment une base,
force $p_0 = p_1 = p_\infty$. Le nom d'aujourd'hui est « automorphisme
préservant l'inertie ».}. Le centralisateur de $\mathfrak{S}_3$ dans
$\mathrm{Out}_{\mathrm{lac}}(\widehat{\pi})$ est noté $\boldsymbol{\Gamma}$ dans
tout le document\footnote{La page 10 le note $\Pi$ (« $\mathrm{Autext}_{\mathrm{lac}}
(\hat\pi_{0,3}, \mathfrak{S}_3) \overset{\text{déf}}{=} \Pi$ »), la page 23
$\boldsymbol{\Gamma}$ ; c'est le même groupe, et $\Pi$ est pris par ailleurs pour
le groupe discret aux pages 15 à 20.}. L'exposant $!$ du dossier
($\hat{\hat{\mathcal{E}}}^{!}$, $\hat{\hat\Gamma}^{!}$) désigne la partie
« pure », où $\sigma = 1$.

\pagerange{1}{3}
\section*{Trois feuillets épars (pages 1 à 3)}

\pagerange{1}{1}
\subsection*{Le groupe $H$ : un commensurateur (page 1)}

Soit $U$ une courbe, $D = \widetilde{U(\mathbb{C})}$ le revêtement universel de
$U(\mathbb{C})$, $\pi$ le groupe des automorphismes de $D$ au-dessus de
$U(\mathbb{C})$, et $G$ un groupe d'automorphismes de $D$ contenant
$\pi$\footnote{La page n'écrit que « $\pi \subset G$ » et ne dit pas ce qu'est
$G$ ; le diagramme qui précède, barré de grandes croix, porte une tour de
revêtements $\overline{U}_{J,\alpha} \to U_{J,\alpha}$ qui n'est pas reprise.}.
La page définit
\[
H = \bigl\{\, u \in G \bigm| \exists\, \pi', \pi'' \subset \pi \text{ d'indice
fini et } v : D/\pi' \xrightarrow{\sim} D/\pi'' \text{ avec } p'' u = v p'
\,\bigr\},
\]
où $p' : D \to D/\pi'$ et $p'' : D \to D/\pi''$ sont les projections :
$u$ \emph{descend} en un isomorphisme entre deux revêtements finis de
$U(\mathbb{C})$. Alors $u \pi' u^{-1} = \pi''$, et l'isomorphisme $v_* : \pi'
\to \pi''$ induit est la conjugaison par $u$. Inversement, si $u \pi u^{-1}
\cap \pi$ est d'indice fini dans $\pi$ et dans $u\pi u^{-1}$, on prend $\pi' =
\pi \cap u^{-1}\pi u$ et $\pi'' = u\pi' u^{-1}$. Donc $H$ est le
\emph{commensurateur} de $\pi$ dans $G$\footnote{Le mot n'est pas sur la page, et
rien n'indique qu'il l'ait eu en tête. Pour $U = \mathbb{P}^1 \setminus \{0, 1,
\infty\}$, $D$ est le demi-plan de Poincaré et $\pi$ est d'indice fini dans
$PSL_2(\mathbb{Z})$ ; dans $G = PSL_2(\mathbb{R})$ son commensurateur est alors
celui de $PSL_2(\mathbb{Z})$, l'image de $GL_2^{+}(\mathbb{Q})$, fait classique
— c'est lui qui rend $\pi$ « arithmétique » au sens du critère de Margulis.}.
Les remarques de la page se lisent ainsi :

\begin{itemize}
\item $\pi'$ et $\pi''$ ont le même indice dans $\pi$ quand $u$ préserve une
mesure pour laquelle $D/\pi$ est de volume fini — par exemple quand $u$ est
une isométrie hyperbolique : $D/\pi'$ et $D/\pi''$ ont alors même aire,
$[\pi:\pi']$ et $[\pi:\pi'']$ fois celle de $D/\pi$\footnote{La page dit
« on aura d'ailleurs $\pi'$ et $\pi''$ de même indice dans $\pi$ », lecture
incertaine, sans hypothèse ; celle-ci est ajoutée. Pour un groupe $G$
quelconque, les deux indices peuvent différer.} ;
\item $v$ fixé, $u$ n'est connu qu'à multiplication à gauche près par un élément
de $\pi''$ (les automorphismes du revêtement $D \to D/\pi''$), et $H$ contient
$\pi$ ;
\item $v$ est déterminé, à homotopie près, par $v_*$ à conjugaison près,
puisque les $D/\pi'$ sont des $K(\pi', 1)$\footnote{La page dit « $v$ est connu
quand on connaît $v_*$ », avec une réserve illisible entre parenthèses (lue
« cas connexe ») ; la précision « à homotopie près » est de cette édition.}.
\end{itemize}

Pour $n \geq 1$, $H_n$ est l'ensemble des $u$ pour lesquels on peut prendre
$\pi'$, $\pi''$ d'indice $n$ ; $H = \bigcup_n H_n$ et $H_n = H_n^{-1}$. La
page affirme que $H_n$ est un sous-groupe et s'arrête sur « et $\pi$
d'indice » ; ce n'est pas évident — pour un produit $uw$ on ne dispose a
priori que de sous-groupes d'indice au plus $n^2$ — et nous ne
l'affirmons pas\footnote{Le paragraphe central (« $H/\pi$ … s'identifie à un
sous-\ldots{} de l'un des “groupoïdes” d'isom. \ldots ») ne se lit que par
fragments ; on n'en tire rien.}.

\pagerange{2}{2}
\subsection*{Automorphismes à lacets commutant à $\sigma_0$ (page 2)}

La page cherche les automorphismes à lacets $u$ de $\widehat{\pi}$ qui
commutent exactement à $\sigma_0$ et, à un automorphisme intérieur près, à
$\rho$. Commuter à $\sigma_0$, qui échange $l_1$ et $l_\infty$, force la forme
\[
u(l_1) = \mathrm{int}(g_1)(l_1^{p}), \qquad
u(l_\infty) = \mathrm{int}\bigl(\sigma_0(g_1)\bigr)(l_\infty^{p}),
\]
avec $p \in \widehat{\mathbb{Z}}^{*}$ et $g_1 \in \widehat{\pi}$ défini modulo
$L_1$ à droite. Posant $h_1 = \sigma_0(g_1)^{-1} g_1$, la page énonce deux
conditions encadrées, où $\beta, \gamma \in \widehat{\mathbb{Z}}$ sont « uniquement
déterminés par $g_1, p$ » :
\[
\begin{aligned}
(1)\quad & l_\infty^{p} \cdot \mathrm{int}(h_1)(l_1^{p}) \cdot
\mathrm{int}\bigl(h_1 l_1^{-\gamma} \rho^{-1}(h_1)\bigr)(l_0^{p}) = 1, \\
(2)\quad & \rho(h_1)\, l_\infty^{-\gamma}\, h_1\, l_1^{-\gamma}\, \rho^{-1}(h_1)
= l_0^{\beta}.
\end{aligned}
\]
La feuille ne fixe ni l'ordre de la relation fondamentale ni la normalisation
qui définit $\gamma$, et ces deux relations ne sont pas vérifiées
ici\footnote{Une note marginale en oblique donne aussi
$u(l_0) = \mathrm{int}\bigl(g_1 l_1 \rho^{-1}(h_1)\bigr)(l_0^{p})$ ; deux autres
ne se lisent que par fragments. L'indice de $u(l_1)$ est surchargé.}. Ce qui
suit en découle, et l'algèbre en est juste : (2) s'écrit $h_1 l_1^{-\gamma}
\rho^{-1}(h_1) = l_\infty^{\gamma} \rho(h_1)^{-1} l_0^{\beta}$, et comme
$l_0^{\beta}$ commute à $l_0^{p}$, (1) devient
\[
(1')\qquad l_\infty^{p} \cdot \mathrm{int}(h_1)(l_1^{p}) \cdot
\mathrm{int}\bigl(l_\infty^{\gamma} \rho(h_1)^{-1}\bigr)(l_0^{p}) = 1 .
\]
Enfin, changer $g_1$ en $g_1 l_1^{\mu}$ change $h_1$ en $h'_1 = l_\infty^{-\mu}
h_1 l_1^{\mu}$, puisque $\sigma_0(l_1) = l_\infty$ ; la page en conclut que
$\gamma$ se décale et que $\beta$ ne change pas\footnote{« (i.e. $\gamma' =
\gamma \ldots$ » se termine par un signe gratté ; la dernière formule de la
page a ses exposants surchargés.}.

\pagerange{3}{3}
\subsection*{Automorphismes compatibles à un morphisme (page 3, barrée)}

La page entière est barrée de deux longs traits. Elle place côte à côte les
espaces classifiants $B_{\pi,\Gamma} \leftarrow B_\pi$ d'une extension et
ceux d'un quotient $\varphi : \pi \to \pi'$, et pose la question suivante : si
$u \in \mathrm{Aut}(\pi)$ est « compatible avec $\varphi$ », existe-t-il un unique
$u' \in \mathrm{Aut}(\pi')$ tel que $u' \varphi = \varphi u$ ? L'unicité est
vraie dès que $\varphi$ est surjectif, et l'existence équivaut alors à
$u(\ker \varphi) = \ker \varphi$\footnote{La page n'énonce pas d'hypothèse sur
$\varphi$, et passe par des essais barrés ($u' \varphi u^{-1} =
\mathrm{int}(g)\,\varphi u$, etc.) qui envisagent la version à automorphisme
intérieur près. Les deux suites exactes de groupes « doublement complétés » qui
terminent la page, $1 \to \pi \to \hat{\hat{\mathfrak{S}}}(\pi, I') \to$ et $1 \to
\pi' \to \hat{\hat{\mathcal{E}}}(\pi') \to$, portent des lettres rondes que la
transcription ne distingue pas sûrement ; elles annoncent, sans le définir, le
$\hat{\hat{\mathcal{E}}}$ des pages 22 et suivantes.}. Le sous-groupe
$\mathcal{E}_\varphi$ de ces $u$ est défini et la page s'arrête.

\pagerange{4}{14}
\section*{Première rédaction : automorphismes extérieurs à lacets de
$\widehat{\pi}_{0,3}$ (pages 4 à 14)}

\pagerange{4}{6}
\subsection*{L'action extérieure de $\Gamma_{0,3}$ (pages 4 à 6)}

Dans toute cette rédaction, $\widehat{\pi}$ est présenté par $l_0, l_1,
l_\infty$ et la relation
\[
l_0\, l_1\, l_\infty = 1 .
\]
Un endomorphisme est donné par trois éléments $l'_0, l'_1, l'_\infty$
satisfaisant la même relation, et c'est un automorphisme dès qu'ils engendrent
topologiquement $\widehat{\pi}$ : un groupe profini de type fini est hopfien,
tout endomorphisme surjectif est bijectif\footnote{La page dit « sauf erreur ce
sera un automorphisme s'ils engendrent $\pi$ (topologiquement) » ; c'est exact,
et la raison est ajoutée.}.

La permutation circulaire $\rho(l_0) = l_1$, $\rho(l_1) = l_\infty$,
$\rho(l_\infty) = l_0$ est un automorphisme d'ordre $3$. La transposition
$\sigma_\infty$ est définie par
\[
\sigma_\infty(l_0) = l_1, \qquad \sigma_\infty(l_1) = l_0, \qquad
\sigma_\infty(l_\infty) = l_0^{-1} l_1^{-1} = \mathrm{int}(l_0^{-1})(l_\infty),
\]
et $\sigma_\infty^2 = \mathrm{id}$ ; on pose $\sigma_0 = \rho\,\sigma_\infty\,
\rho^{-1}$ et $\sigma_1 = \rho\,\sigma_0\,\rho^{-1}$, de sorte que $\rho
\sigma_i \rho^{-1} = \sigma_{\rho(i)}$. Explicitement, $\sigma_0$ échange $l_1$
et $l_\infty$ et envoie $l_0$ sur $l_1^{-1} l_\infty^{-1}$. Le sous-groupe
$\mathfrak{S}_3^{+} \simeq \mathbb{Z}/3$ agit donc vraiment. Mais pas
$\mathfrak{S}_3$ : dans $\mathfrak{S}_3$ on a $\sigma_0\sigma_1 =
\sigma_1\sigma_\infty = \sigma_\infty\sigma_0 = \rho$, tandis que dans
$\mathrm{Aut}(\widehat{\pi})$
\[
\sigma_0\sigma_1 = \mathrm{int}(l_1^{-1})\,\rho, \qquad
\sigma_1\sigma_\infty = \mathrm{int}(l_\infty^{-1})\,\rho, \qquad
\sigma_\infty\sigma_0 = \mathrm{int}(l_0^{-1})\,\rho .
\]
$\mathfrak{S}_3$ n'agit que par automorphismes extérieurs. Pour le groupe
discret, on ne peut pas mieux faire : une action vraie de $\mathfrak{S}_3$ sur
$\pi$ relevant cette action extérieure scinderait l'extension correspondante,
et les sous-groupes finis de celle-ci fixent un point du revêtement universel ;
or $\mathfrak{S}_3$ n'a pas de point fixe sur $\mathbb{P}^1_{\mathbb{C}}
\setminus \{0, 1, \infty\}$\footnote{L'argument par les sous-groupes finis est
de cette édition ; la page invoque seulement l'absence de point fixe.}. Si
$\mathbb{Z}/3$ agit vraiment, c'est que son action sur la sphère a deux points
fixes, $-j$ et $-j^2 = -\bar j$, les racines primitives sixièmes de
l'unité $e^{\pm i\pi/3}$, et qu'on peut prendre l'un d'eux pour point
base\footnote{La page écrit « savoir $\pm j$ (racine primitive troisième de
1) ». Mais $\rho(z) = 1/(1-z)$ a pour points fixes les racines de $z^2 - z + 1
= 0$, soit $e^{\pm i\pi/3} = -j^2, -j$ ; $j$ lui-même n'est pas fixe.}.

La conjugaison complexe, élément du groupe de Teichmüller $\Gamma_{0,3}$,
agit par
\[
\tau(l_0) = l_0^{-1}, \qquad \tau(l_1) = l_1^{-1}, \qquad \tau(l_\infty) =
l_1 l_0 = \mathrm{int}(l_1)(l_\infty^{-1}),
\]
automorphisme à lacets de multiplicateur $-1$, avec $\tau^2 = 1$. Il commute à
$\mathfrak{S}_3$ modulo automorphismes intérieurs, d'où une action extérieure
de $\mathfrak{S}_3 \times \{\pm 1\} = \Gamma_{0,3}$, qui est l'action
canonique\footnote{« Groupe de Teichmüller » est le nom qu'il donne à
$\Gamma_{0,3}$ ; c'est le groupe modulaire étendu (« extended mapping class
group ») de la sphère à trois trous. Que ce soit tout
$\mathrm{Out}_{\mathrm{lac}}(\pi)$ pour le groupe discret est le théorème de
Dehn--Nielsen--Baer.}.

\pagerange{6}{11}
\subsection*{Les automorphismes extérieurs qui commutent à $\mathfrak{S}_3$
(pages 6 à 11)}

\emph{Normalisation.} Tout automorphisme extérieur à lacets de multiplicateur
$p$ et de permutation triviale sur $0$ se représente par un $u$ tel que
\[
(*)\qquad u(l_0) = l_0^{p},
\]
et ce représentant est unique à composition près par $\mathrm{int}(\lambda)$,
$\lambda \in L_0$, puisque le normalisateur de $L_0$ dans $\widehat{\pi}$ est
$L_0$\footnote{La page dit que $u$ « normalise $L_0$ », puis « la classe de
conj. de $l_0$ », et affirme l'unicité sans la justifier. Le fait utilisé — un
sous-groupe d'inertie d'un groupe libre profini est son propre normalisateur —
est classique.}.

\emph{Commutation à $\rho$.} Si la classe de $u$ commute à celle de $\rho$, il
existe un unique $g \in \widehat{\pi}$ tel que
\[
(**)\qquad u\,\rho = \mathrm{int}(g)\,\rho\, u ,
\]
et remplacer $u$ par $\mathrm{int}(\lambda) u$, $\lambda = l_0^{\alpha}$,
remplace $g$ par
\[
\lambda\, g\, \rho(\lambda)^{-1} = l_0^{\alpha}\, g\, l_1^{-\alpha},
\qquad \alpha \in \widehat{\mathbb{Z}}.
\]
L'exposant $\alpha$ parcourt $\widehat{\mathbb{Z}}$, puisque $\lambda$
parcourt $L_0$\footnote{La page écrit $\alpha \in \widehat{\mathbb{Z}}^{*}$.}.
Réciproquement, $(*)$ et $(**)$ déterminent $u$ en fonction de $p$ et $g$ :
\[
u(l_0) = l_0^{p}, \qquad u(l_1) = \mathrm{int}(g)(l_1^{p}), \qquad u(l_\infty)
= \mathrm{int}\bigl(g\rho(g)\bigr)(l_\infty^{p}),
\]
car $u(l_1) = u\rho(l_0) = \mathrm{int}(g)\rho(l_0^{p})$ et $u(l_\infty) =
u\rho(l_1) = \mathrm{int}(g)\,\mathrm{int}(\rho(g))\,\rho(l_1^{p})$ ; $(**)$ est
alors vraie, les deux membres coïncidant sur $l_0$ et $l_1$. Pour que ces
formules définissent un endomorphisme, il faut et il suffit de la relation
\[
\boxed{\; l_0^{p} \cdot \mathrm{int}(g)(l_1^{p}) \cdot
\mathrm{int}\bigl(g\rho(g)\bigr)(l_\infty^{p}) = 1 \;}
\]
soit, en développant puis en conjuguant par $g^{-1}$,
$\mathrm{int}(g^{-1})(l_0^{p})\cdot l_1^{p} \cdot
\mathrm{int}(\rho(g))(l_\infty^{p}) = 1$. Pour un automorphisme, il faut de plus
que les trois images engendrent $\widehat{\pi}$ ; « Ça semble délicat… », dit
la page, et le dossier n'y revient pas\footnote{Il n'y a en effet pas de
critère maniable de génération topologique pour un groupe libre profini ;
c'est une des raisons pour lesquelles la théorie ultérieure travaille avec des
paires $(\lambda, f)$ dont on \emph{impose} l'inversibilité.}.

Pour $\tau$, $p = -1$ et $g = l_1$ : on vérifie $\tau\rho = \mathrm{int}(l_1)\,
\rho\,\tau$.

\emph{Commutation à $\sigma_\infty$.} $\mathfrak{S}_3$ étant engendré par $\rho$
et $\sigma_\infty$, reste à écrire qu'il existe $h \in \widehat{\pi}$ avec
$u\,\sigma_\infty = \mathrm{int}(h)\,\sigma_\infty\, u$, et il suffit de
l'évaluer en $l_0$ et $l_1$. En $l_0$ : $\mathrm{int}(g)(l_1^{p}) =
\mathrm{int}(h)(l_1^{p})$, donc $g^{-1}h$ centralise $l_1^{p}$, donc est dans
$L_1$ : $h = g\, l_1^{q}$ avec $q \in \widehat{\mathbb{Z}}$. En $l_1$ :
$l_0^{p} = \mathrm{int}\bigl(h\,\sigma_\infty(g)\bigr)(l_0^{p})$, d'où la
condition
\[
\boxed{\; g\, l_1^{q}\, \sigma_\infty(g) = l_0^{r} \;} \qquad (q, r \in
\widehat{\mathbb{Z}}),
\]
qui ne dépend que de $g$, pas de $p$\footnote{La page 9 écrit « $\exists\, q, r
\in \hat{\mathbb{Z}}^{*}$ » ; la page 8 avait correctement $q \in
\hat{\mathbb{Z}}$. Les deux exposants sont dans $\widehat{\mathbb{Z}}$ ; pour
$\tau$ ce sont des unités, mais rien ne l'impose en général.}. Le changement
$g \mapsto l_0^{\alpha} g l_1^{-\alpha}$ la laisse invariante avec les mêmes
$(q, r)$ — ils ne dépendent que de l'automorphisme extérieur. Pour $\tau$, $l_1
l_1^{q} l_0 = l_0^{r}$ donne $q = -1$, $r = 1$.

La note marginale de la page 9 propose d'écrire plutôt la relation sous la
forme $\sigma_\infty(g^{-1}) = l_0^{\beta} g\, l_1^{\gamma}$. Tirant
$\sigma_\infty(g) = l_1^{-q} g^{-1} l_0^{r}$ et inversant, on trouve $\beta =
-r$ et $\gamma = q$\footnote{La marge écrit « i.e. $\beta = r$, $\gamma = -q$ » :
les deux signes sont à changer. La page 10 donne d'ailleurs, pour $\tau$,
« $\beta = -1$ », ce qui est $-r$.}. Mieux : $\beta = \gamma$. En appliquant
l'involution $\sigma_\infty$ à la relation, on obtient $g\, l_1^{\beta-\gamma}
g^{-1} = l_0^{\beta-\gamma}$, et dans l'abélianisé, où $\bar l_0$ et $\bar l_1$
sont indépendants, cela force $\beta = \gamma$. La relation de commutation à
$\mathfrak{S}_3$ est donc
\[
\sigma_\infty(g^{-1}) = l_0^{\beta}\, g\, l_1^{\beta},
\]
exactement comme l'écrit la page 11, et le couple $(q, r)$ de la page 9 est
un seul invariant : $r = -q = -\beta$\footnote{La page 11 conclut « Ce qui
détermine $\beta$ de façon unique en fonction de $g$ » ; l'égalité $\beta =
\gamma$ n'y est pas justifiée, et la réduction de $(q, r)$ à un seul nombre
n'est pas dite. Pour $\tau$ : $\beta = -1$, et en effet $\sigma_\infty(l_1^{-1})
= l_0^{-1} = l_0^{-1}\, l_1\, l_1^{-1}$.}.

\emph{Le cas discret.} Pour le groupe discret $\pi$, les seules solutions sont
l'identité et $\tau$ : le groupe des automorphismes extérieurs à lacets de
$\pi$ est $\Gamma_{0,3}$ (Dehn--Nielsen--Baer), et le centralisateur de
$\mathfrak{S}_3$ dans $\mathfrak{S}_3 \times \mathbb{Z}/2$ est $\{1, \tau\}$.
La page décrit la classe de $g$ pour $\tau$ : $g = l_0^{a} l_1^{b}$ avec $a + b
= 1$, orbite de $l_1$ sous $g \mapsto l_0^{a} g l_1^{-a}$\footnote{La page note
$\alpha, \beta$ ces deux exposants, ce qui télescope le $\beta$ de la
relation ; on les renomme. La fin de la seconde ligne du système (« au moins,
$\beta = -1$ ») se lit mal.}.

\emph{Le groupe $\boldsymbol{\Gamma}$.} Normaliser en $L_1$ plutôt qu'en $L_0$
revient à remplacer $u$ par $\mathrm{int}(g^{-1})u$, et $g$ par $\rho(g)$ :
on vérifie $\mathrm{int}(g^{-1})u\,\rho = \mathrm{int}(\rho(g))\,\rho\,
\mathrm{int}(g^{-1})u$. Sur le groupe $\boldsymbol{\Gamma}$ des automorphismes
extérieurs à lacets commutant à $\mathfrak{S}_3$ on a donc le multiplicateur
$\chi : \boldsymbol{\Gamma} \to \widehat{\mathbb{Z}}^{*}$, l'invariant $\beta$ à
valeurs dans $\widehat{\mathbb{Z}}$, et la classe $\gamma(u)$ de $g$ dans
l'ensemble des orbites de $\widehat{\mathbb{Z}}$ agissant sur $\widehat{\pi}$
par $\varphi(\alpha)\cdot g = l_0^{\alpha} g\, l_1^{-\alpha}$. L'application
\[
(\gamma, \chi) : \boldsymbol{\Gamma} \longrightarrow
\widehat{\pi}/\varphi(\widehat{\mathbb{Z}}) \times \widehat{\mathbb{Z}}^{*}
\]
est injective, puisque $(p, g)$ détermine $u$, et son image est formée des
$(g, p)$ satisfaisant les deux relations encadrées (et la condition de
génération)\footnote{La page écrit le but $\pi/\hat{\mathbb{Z}} \times
\hat{\mathbb{Z}}^{*}$ ; ce n'est pas un quotient de groupe mais un ensemble
d'orbites. La note marginale de la page 10 (« $\gamma : E_0 \to \pi$ \ldots{}
cocycle ») ne se lit que par fragments ; la structure de cocycle est reprise
proprement aux pages 24 et 36.}.

On reconnaît là, vingt ans avant, la forme des deux premières relations de
définition du groupe de Grothendieck--Teichmüller $\widehat{GT}$ : dans la
présentation de Drinfeld (1990), un élément $(\lambda, f)$ agit par $x \mapsto
x^{\lambda}$, $y \mapsto f^{-1} y^{\lambda} f$, et ses relations (I) et (II)
traduisent la commutation, modulo automorphismes intérieurs, à la transposition
$x \leftrightarrow y$ et à la permutation circulaire\footnote{Le rapprochement
est de cette édition. La normalisation diffère : ici $g$ est fixé par $(**)$,
non par la condition « $f$ dans le groupe dérivé » de Drinfeld, de sorte que
les équations ne se superposent pas terme à terme. La troisième relation de
$\widehat{GT}$, le pentagone, porte sur $\pi_{0,5}$ et n'a pas de trace dans
ce dossier. Grothendieck ne pouvait connaître ni Drinfeld (1990) ni Ihara
(1991).}.

\pagerange{11}{14}
\subsection*{Les racines quasi-carrées $\varepsilon_i$ (pages 11 à 14)}

Dans l'extension de $\mathfrak{S}_3$ par $\widehat{\pi}$ que forment, dans
$\mathrm{Aut}(\widehat{\pi})$, les $\mathrm{int}(g)$ et les $\sigma_i$, la page
cherche de « plus jolis générateurs » que les six $l_i, \sigma_i$ : pour chaque
$i$, un élément $\varepsilon_i$ au-dessus de $\sigma_i$, qui fixe $l_i$ et dont le
carré est $l_i$. Puisque $\mathrm{int}(l_0)\sigma_\infty$ fixe $l_\infty$, on
cherche $\varepsilon_\infty = l_0\,\sigma_\infty\, l_\infty^{\alpha}$, et
$\varepsilon_\infty^2 = l_\infty$ impose $\alpha = 1$\footnote{La page passe par
plusieurs essais barrés, dont un « prenons $\alpha = 1$ ? » en marge, un calcul
refusé (« Non ! ») et un calcul accepté (« ok »). La note « doit écrire
$\varepsilon_0^2 = l_\infty$ » a son indice mal lu : la relation cherchée est
$\varepsilon_\infty^2 = l_\infty$.}. D'où, par $\rho$,
\[
\varepsilon_\infty = l_0\,\sigma_\infty\, l_\infty, \qquad
\varepsilon_0 = l_1\,\sigma_0\, l_0, \qquad
\varepsilon_1 = l_\infty\,\sigma_1\, l_1, \qquad \varepsilon_i^2 = l_i,
\]
et $\rho\,\varepsilon_i\,\rho^{-1} = \varepsilon_{\rho(i)}$. On a en retour
$\sigma_0 = \varepsilon_0 \varepsilon_1^2$, $\sigma_1 = \varepsilon_1
\varepsilon_\infty^2$, $\sigma_\infty = \varepsilon_\infty \varepsilon_0^2$, et
les relations
\[
\varepsilon_0^2\,\varepsilon_1^2\,\varepsilon_\infty^2 = 1, \qquad
(\varepsilon_0\varepsilon_1^2)^2 = (\varepsilon_1\varepsilon_\infty^2)^2 =
(\varepsilon_\infty\varepsilon_0^2)^2 = 1 .
\]
De $\rho = l_1\sigma_0\sigma_1$ on tire $\rho = \varepsilon_1^2\,\varepsilon_0\,
\varepsilon_1\, \varepsilon_0^{-2}$, et ses deux conjugués circulaires. En
multipliant les trois expressions et en écrivant $\rho^3 = 1$, on trouve
\[
(\varepsilon_0\,\varepsilon_1\,\varepsilon_\infty)^2 = 1 .
\]

La page juge cela « inquiétant, car il n'y a pas d'élément d'ordre fini dans le
groupe discret », en conclut $\varepsilon_0\varepsilon_1\varepsilon_\infty = 1$,
puis $\sigma_\infty\sigma_1\sigma_0 = 1$, « relation inattendue » — et barre
tout ce passage. Il a raison de le barrer : les deux relations sont fausses. Le
produit $\varepsilon_0\varepsilon_1\varepsilon_\infty$ n'est pas dans
$\widehat{\pi}$, il est au-dessus de la transposition $\sigma_0\sigma_1
\sigma_\infty = \sigma_1$, et une extension de $\mathfrak{S}_3$ a bien des
éléments d'ordre $2$ ; en fait
\[
\varepsilon_0\,\varepsilon_1\,\varepsilon_\infty = \sigma_1 ,
\]
comme le dit la ligne non barrée qui suit (« Sans doute
$\varepsilon_0\varepsilon_1\varepsilon_\infty = \sigma_1$ »), et
$\sigma_\infty\sigma_1\sigma_0$ est une transposition, non l'identité. Avec
$\sigma_1 = \varepsilon_1\varepsilon_\infty^2$, cette égalité équivaut à
$\varepsilon_0 = \varepsilon_1\varepsilon_\infty\varepsilon_1^{-1}$, et l'on a
de même
\[
\varepsilon_0 = \varepsilon_1\,\varepsilon_\infty\,\varepsilon_1^{-1}, \qquad
\varepsilon_1 = \varepsilon_\infty\,\varepsilon_0\,\varepsilon_\infty^{-1},
\qquad \varepsilon_\infty = \varepsilon_0\,\varepsilon_1\,\varepsilon_0^{-1},
\]
les trois relations que la page écrit à sa dernière ligne avant la
critique\footnote{Le second membre de « Sans doute » est un signe court, que la
transcription lit $\sigma_1$ « sans conviction » ; le calcul confirme cette
lecture. Toutes les identités de cette section ont été vérifiées dans le
groupe libre de rang $2$, dans la convention $l_0 l_1 l_\infty = 1$ de cette
rédaction.}.

\pagerange{14}{22}
\section*{Seconde rédaction : l'extension discrète $\mathcal{E}_{0,3}$ (pages 14
à 22)}

\pagerange{14}{16}
\subsection*{La critique, et la géométrie de la sphère (pages 14 à 16)}

« Critique des notations et de l'approche » : les éléments $l_i$, $\sigma_i$,
$\varepsilon_i$, $\rho$ vivent dans le groupe discret, qu'il faut distinguer du
groupe profini, et pour le groupe discret — et, dans une certaine mesure, pour
le profini — il faut utiliser le langage géométrique du groupoïde fondamental.

On munit $\mathbb{P}^1(\mathbb{C})$ de sa structure de sphère orientée pour
laquelle $\rho$ est la rotation d'angle $+2\pi/3$ autour de l'axe $P\overline{P}$,
où $P, \overline{P} = e^{\pm i\pi/3}$ sont ses deux points fixes ; les
$\sigma_i$ sont les demi-tours autour des axes $(0, 2)$, $(1, -1)$ et
$(\infty, \frac12)$, et $\tau$ la symétrie par rapport au plan du cercle réel.
On pose $U = U_{0,3} = \mathbb{P}^1(\mathbb{C}) \setminus \{0, 1, \infty\}$ et
$\pi = \pi_1(U, P)$.

\emph{Le groupe fondamental mixte.} L'extension $\mathcal{E} =
\pi_1(U, \Gamma; P)$ de $\Gamma = \mathfrak{S}_3 \times \mathbb{Z}/2$ par $\pi$
est définie « à la Malgoire » : ses éléments sont les couples $(\gamma, l)$ d'un
$\gamma \in \Gamma$ et d'une classe d'homotopie de chemins $l$ de $P$ à
$\gamma^{-1}(P)$, avec
\[
(\gamma, l)(\gamma', l') = \bigl(\gamma\gamma',\; \gamma'^{-1}(l) \circ l'\bigr),
\]
les chemins se composant comme des fonctions\footnote{La page note d'abord la
classe $(\gamma, \lambda)$, puis $(\gamma, l)$. La convention « comme des
fonctions » est celle qui fait opérer le groupoïde fondamental à gauche sur
les fibres d'un revêtement. J. Malgoire est son collègue de Montpellier ; le
dossier ne dit pas à quel texte il renvoie — vraisemblablement aux
\emph{Cartes cellulaires} de Malgoire et C. Voisin (1977).}. C'est ce qu'on
appelle aujourd'hui le groupe fondamental orbifold (ou équivariant) de $U$ sous
$\Gamma$ : le $\pi_1$ de la construction de Borel $U_{h\Gamma}$, qui s'insère
dans $1 \to \pi \to \mathcal{E} \to \Gamma \to 1$. Comme $\pi$ est de centre
trivial et que le centralisateur de $\pi$ dans $\mathcal{E}$ est trivial,
$\mathcal{E}$ se plonge dans $\mathrm{Aut}(\pi)$, et c'est là qu'on vérifie les
relations ci-dessous.

\emph{Les lacets.} On relève les $\sigma_i$ en des éléments de $\mathcal{E}$ en
les accompagnant du demi-grand cercle de $P$ à $\overline{P} = \sigma_i(P)$
fixé par la symétrie $\tau\sigma_i$ ; $\rho$ se relève avec le chemin constant.
Soit $l_i$ le lacet de $P$ qui va vers $i$ par un arc de grand cercle, en fait
le tour dans le sens direct et revient. Alors
\[
(1)\qquad l_\infty\, l_1\, l_0 = 1,
\]
« et non $l_0 l_1 l_\infty = 1$ ! » : la composition « comme des fonctions »
renverse l'ordre de la première rédaction.

\pagerange{17}{19}
\subsection*{Présentation de $\mathcal{E}^{+}$, et les $\varepsilon_i$ revus
(pages 17 à 19)}

On note $\mathcal{E}^{+}$ l'image inverse de $\mathfrak{S}_3$. Les relevés
satisfont
\[
(2)\qquad \sigma_0^2 = \sigma_1^2 = \sigma_\infty^2 = 1, \qquad
\sigma_0\sigma_1 = \rho\, l_0, \quad \sigma_1\sigma_\infty = \rho\, l_1, \quad
\sigma_\infty\sigma_0 = \rho\, l_\infty,
\]
avec $\rho\sigma_i\rho^{-1} = \sigma_{\rho(i)}$, et les conjugaisons
\[
(7)\qquad \mathrm{int}(\rho)(l_i) = l_{\rho(i)}, \qquad
\mathrm{int}(\sigma_0) : l_1 \leftrightarrow l_\infty,\;
l_0 \mapsto l_\infty^{-1} l_1^{-1} = \mathrm{int}(l_1)(l_0),
\]
et de même, par $\rho$, pour $\sigma_1$ et $\sigma_\infty$\footnote{Les seconds
membres de (2) sont récrits sur la page : on y devine une première écriture
$l_0\rho$, $l_1\rho$, $l_\infty\rho$, corrigée en $\rho\, l_0$, $\rho\, l_1$,
$\rho\, l_\infty$. La correction est juste ; la première écriture reviendra à
la page 29. La note marginale de la page 17 donne $\sigma_0(l_0) = l_0^{-1}
l_1^{-1}$, qui n'est pas compatible avec (1) ; celle de la page 19 donne la
bonne valeur $l_\infty^{-1} l_1^{-1}$. La formule qui introduit (7) (« comme
$\sigma_i = \ldots$ ») est mal lue.}.

On ne voit pas d'expression de $\rho$ par les seuls $\sigma_i$, et pour cause :
une première affirmation, que les relations d'une présentation de
$\mathfrak{S}_3$ jointes à (1) et (2) présentent $\mathcal{E}^{+}$, est barrée et
jugée « fausse » en marge. La page 19 propose alors : générateurs $\sigma_0,
\sigma_1, \sigma_\infty, \rho$, et relations
\[
(4)\quad (\rho^{-1}\sigma_\infty\sigma_0)(\rho^{-1}\sigma_1\sigma_\infty)
(\rho^{-1}\sigma_0\sigma_1) = 1, \qquad
(5)\quad \sigma_0^2 = \sigma_1^2 = \sigma_\infty^2 = 1, \qquad
(6)\quad \rho^3 = 1,
\]
où (4) est (1) après élimination des $l_i$ par (2). Ces relations sont vraies
dans $\mathcal{E}^{+}$, mais elles ne le présentent pas : le groupe qu'elles
définissent a pour abélianisé $(\mathbb{Z}/2)^3 \times \mathbb{Z}/3$, alors que
celui de $\mathcal{E}^{+}$ est $\mathbb{Z}/6$. Il faut ajouter les relations
$\rho\sigma_i\rho^{-1} = \sigma_{\rho(i)}$ — ce que dit la note marginale de la
même page (« Il faut ajouter la relation \ldots{} $\rho\sigma_0\rho^{-1} =
\sigma_1$ »). Elles rendent $\sigma_1, \sigma_\infty$ superflus, et (4) devient
alors conséquence de $\sigma_0^2 = \rho^3 = 1$ ; il reste
\[
\mathcal{E}^{+} = \langle\, \sigma_0, \rho \mid \sigma_0^2 = \rho^3 = 1 \,\rangle
\simeq \mathbb{Z}/2 * \mathbb{Z}/3 \simeq PSL_2(\mathbb{Z}) .
\]
En effet, la fonction modulaire $\lambda$ identifie $U$ au quotient du
demi-plan de Poincaré par l'image $\overline{\Gamma(2)}$ du sous-groupe de
congruence $\Gamma(2)$, et l'action anharmonique de $\mathfrak{S}_3$ sur $U$ à
celle de $PSL_2(\mathbb{Z})/\overline{\Gamma(2)} \simeq PSL_2(\mathbb{F}_2)$ ;
le groupe des relèvements au demi-plan, qui est $\mathcal{E}^{+}$, est donc
$PSL_2(\mathbb{Z})$. Celui-ci est engendré par $\sigma_0$ et $\rho$, ce qui
donne une surjection de $\mathbb{Z}/2 * \mathbb{Z}/3$ sur un groupe qui lui est
isomorphe ; comme il est de type fini et résiduellement fini, donc hopfien,
c'est un isomorphisme\footnote{Cette identification, et l'argument, sont de
cette édition ; le dossier ne nomme pas $PSL_2(\mathbb{Z})$. La remarque de la
page 19 — si l'on fait $\sigma_i = 1$ il ne reste que $\rho^3 = 1$, « ce qui
prouve que $\rho$ n'est pas » exprimable par les $\sigma_i$ — est juste, et
vaut aussi pour la présentation complétée : le quotient par le sous-groupe
distingué engendré par $\sigma_0$ est $\mathbb{Z}/3$.}.

\emph{Les $\varepsilon_i$, seconde version (page 18).} Dans les nouvelles
conventions,
\[
(8)\qquad \varepsilon_0 = l_0\,\sigma_0\, l_1, \qquad
\varepsilon_1 = l_1\,\sigma_1\, l_\infty, \qquad
\varepsilon_\infty = l_\infty\,\sigma_\infty\, l_0,
\]
avec $\varepsilon_i^2 = l_i$ (la page le vérifie pour $\varepsilon_0$, par
$\sigma_0(l_1 l_0) = l_1^{-1}$) et $\rho\varepsilon_i\rho^{-1} =
\varepsilon_{\rho(i)}$\footnote{Dans le premier (8), barré, de la page 19 comme
dans celui de la page 18, les exposants des $l$ sont surchargés et illisibles ;
la vérification « on a bien $\varepsilon_0^2 = l_0$ » de la page 18 fixe les
exposants à $1$. La page note que « les $\varepsilon_i$ ici sont les
$\varepsilon_i^{-1}$ de tantôt », ce qui tient au renversement de l'ordre des
lacets entre les deux rédactions.}. On a aussi $\varepsilon_1 =
\sigma_0\,\rho$ : les $\varepsilon_i$ sont, dans $PSL_2(\mathbb{Z})$, les
générateurs paraboliques des stabilisateurs de pointes, dont les carrés
engendrent ceux de $\overline{\Gamma(2)}$ — c'est pourquoi ils sont « racines
carrées » des lacets. Le calcul de la page donne
\[
(9)\qquad
\varepsilon_0\varepsilon_1 = l_0\, l_\infty^2\, \rho\, l_0\, l_\infty, \qquad
\varepsilon_1\varepsilon_\infty = l_1\, l_0^2\, \rho\, l_1\, l_0, \qquad
\varepsilon_\infty\varepsilon_0 = l_\infty\, l_1^2\, \rho\, l_\infty\, l_1,
\]
et, en éliminant les $l_i = \varepsilon_i^2$, la page 21 écrit le système
\[
(10)\qquad \rho^3 = 1, \quad
\varepsilon_\infty^2\varepsilon_1^2\varepsilon_0^2 = 1, \quad
\begin{aligned}
\varepsilon_1 &= \varepsilon_0\,\varepsilon_\infty^{4}\,\rho\,\varepsilon_0^{2}\,\varepsilon_\infty^{2},\\
\varepsilon_\infty &= \varepsilon_1\,\varepsilon_0^{4}\,\rho\,\varepsilon_1^{2}\,\varepsilon_0^{2},\\
\varepsilon_0 &= \varepsilon_\infty\,\varepsilon_1^{4}\,\rho\,\varepsilon_\infty^{2}\,\varepsilon_1^{2},
\end{aligned}
\]
qui est exactement (1) et (9) récrits\footnote{Le système (10) est lu « tel quel
et très incertain » ; les exposants lus $4$ sont ceux que (9) impose. À la
deuxième des trois dernières lignes, la transcription lit (avec doute)
$\varepsilon_\infty^2$ au pénultième facteur ; (9) exige $\varepsilon_1^2$, et
c'est ce qui est écrit ici. Les trois lignes ont été vérifiées dans le groupe
libre.}. La note marginale de la page 21 demande d'y ajouter « (au moins !) »
les relations $\rho\varepsilon_0\rho^{-1} = \varepsilon_1$, etc. ; avec elles,
$\mathcal{E}^{+} = \langle \varepsilon_1, \rho \mid \rho^3 =
(\varepsilon_1\rho^{-1})^2 = 1 \rangle$, puisque $\sigma_0 =
\varepsilon_1\rho^{-1}$.

\pagerange{20}{21}
\subsection*{Le relèvement canonique du stabilisateur de $P$ (pages 21 et 20)}

Le stabilisateur de $P$ dans $\Gamma$ est d'ordre $6$ : $\{1, \dot\rho,
\dot\rho^2, \dot\sigma_0\tau, \dot\sigma_1\tau, \dot\sigma_\infty\tau\}$, car
les $\sigma_i$ et $\tau$ échangent $P$ et $\overline{P}$. C'est l'image de la
section
\[
(11)\qquad \varphi : \mathfrak{S}_3 \longrightarrow \Gamma, \qquad \sigma
\longmapsto \bigl(\sigma, \tau^{\,(1-\operatorname{sg}\sigma)/2}\bigr),
\]
qu'on note $\Gamma_P$. Puisque $P$ est fixe, ses éléments se relèvent
canoniquement dans $\mathcal{E}$, avec le chemin constant :
\[
(12)\qquad \sigma \longmapsto \tilde\sigma : \mathfrak{S}_3 \longrightarrow
\mathcal{E}, \qquad \tilde\rho = \rho,
\]
et $\tilde\sigma_i$ est la symétrie par rapport au plan $(P, \overline{P},
i)$. On a donc
\[
(13)\qquad \tilde\sigma_0^2 = \tilde\sigma_1^2 = \tilde\sigma_\infty^2 = 1,
\qquad \tilde\sigma_0\tilde\sigma_1 = \tilde\sigma_1\tilde\sigma_\infty =
\tilde\sigma_\infty\tilde\sigma_0 = \rho,
\]
d'où $\rho^3 = 1$, et l'image inverse de $\Gamma_P$ dans $\mathcal{E}$ est le
produit semi-direct $\pi \rtimes \mathfrak{S}_3$\footnote{La page 20 écrit
« l'image inverse $\widetilde\Sigma$ de $\varphi(\mathfrak{S}_3) = \Gamma$ » ;
$\varphi(\mathfrak{S}_3)$ est $\Gamma_P$, d'indice $2$ dans $\Gamma$. La
parenthèse « (ss-groupe diédral 2) » est incertaine.}. Les $\tilde\sigma_i$
renversent l'orientation, donc les lacets :
\[
(14)\qquad \tilde\sigma_0\, l_0\, \tilde\sigma_0^{-1} = l_0^{-1}, \qquad
\tilde\sigma_0\, l_1\, \tilde\sigma_0^{-1} = l_\infty^{-1}, \qquad
\tilde\sigma_0\, l_\infty\, \tilde\sigma_0^{-1} = l_1^{-1},
\]
et de même, par $\rho$, pour $\tilde\sigma_1, \tilde\sigma_\infty$. Les
relations (13), (14) et (1) présentent cette image inverse. Enfin
\[
(15)\qquad \tau_i = \sigma_i\,\tilde\sigma_i = \tilde\sigma_i\,\sigma_i
\qquad (i = 0, 1, \infty), \qquad (16)\qquad \tau_i^2 = 1,
\]
trois relevés de $\tau$ dans $\mathcal{E}$, avec $\rho\,\tau_i\,\rho^{-1} =
\tau_{\rho(i)}$\footnote{Que $\sigma_i$ et $\tilde\sigma_i$ commutent est
affirmé par la page et vérifié ici ; les trois notes marginales en oblique de
la page 20 ne se lisent que par fragments.}.

\pagerange{22}{22}
\subsection*{Le groupe cartographique (page 22)}

L'image inverse $\mathcal{E}_\tau$ de $\{1, \tau\}$ est engendrée par $\tau_0,
\tau_1, \tau_\infty$, avec
\[
(17)\qquad l_0 = \tau_\infty\,\tau_1, \qquad l_1 = \tau_0\,\tau_\infty, \qquad
l_\infty = \tau_1\,\tau_0,
\]
et les seules relations (16) : $\mathcal{E}_\tau \simeq \mathbb{Z}/2 *
\mathbb{Z}/2 * \mathbb{Z}/2$, et $\pi$ est le noyau du morphisme $\tau_i \mapsto
-1$ vers $\{\pm 1\}$, libre sur les $l_i$ définis par (17). C'est le
« groupe cartographique de Malgoire » : le groupe qui décrit les cartes tracées
sur une surface par son action sur les drapeaux, et dont le sous-groupe
d'indice $2$ des mots de longueur paire est $\pi_{0,3}$\footnote{Deux mots du
nom (« groupe cartographique “\ldots{} \ldots” de Malgoire ») sont illisibles.
Le groupe cartographique est exposé dans l'\emph{Esquisse d'un programme}
(1984) ; la relation $l_\infty l_1 l_0 = 1$ découle de (17) et de (16).}.

\pagerange{22}{25}
\section*{Le complété et les automorphismes à lacets (pages 22 à 25)}

Soit $\widehat{\mathcal{E}}$ le complété profini de $\mathcal{E}$, extension de
$\widehat{\Gamma} = \Gamma$ par $\widehat{\pi}$, et $\hat{\hat{\mathcal{E}}} =
\mathrm{Aut}_{\mathrm{lac}}(\widehat{\pi})$, extension de $\hat{\hat\Gamma} =
\mathrm{Out}_{\mathrm{lac}}(\widehat{\pi})$ par $\widehat{\pi}$ :
\[
(18)\qquad 1 \longrightarrow \widehat{\pi} \longrightarrow
\hat{\hat{\mathcal{E}}} \longrightarrow \hat{\hat\Gamma} \longrightarrow 1,
\]
qui contient $\widehat{\mathcal{E}}$ comme sous-extension. Soit
$\boldsymbol{\Gamma}$ le centralisateur de $\mathfrak{S}_3$ dans
$\hat{\hat\Gamma}$ — le groupe des pages 6 à 11 — et $\mathbf{E}$ son image
inverse dans $\hat{\hat{\mathcal{E}}}$. Alors $\Gamma \cap \boldsymbol{\Gamma}
= \{1, \tau\}$ et $\boldsymbol{\Gamma} \cap \mathfrak{S}_3 = \{1\}$, puisque
$\mathfrak{S}_3$ est de centre trivial\footnote{La lecture « $\Gamma \cap
\boldsymbol{\Gamma}$ » est incertaine : le second $\Gamma$, à double jambage,
recouvre un $\mathfrak{S}_3$.}.

Un élément de $\boldsymbol{\Gamma}$ agit trivialement sur les trois classes
d'inertie (sa permutation commute à tout $\mathfrak{S}_3$), donc
$\boldsymbol{\Gamma} \subset \hat{\hat\Gamma}^{!}$, la partie pure ; et comme
$\mathfrak{S}_3$ se relève dans $\hat{\hat\Gamma}$, $\hat{\hat\Gamma} =
\hat{\hat\Gamma}^{!} \rtimes \mathfrak{S}_3 \supset \mathfrak{S}_3 \times
\boldsymbol{\Gamma}$. La note marginale de la page 23 dit : « j'ignore si
elles coïncident », c'est-à-dire si $\boldsymbol{\Gamma} =
\hat{\hat\Gamma}^{!}$ — si tout automorphisme extérieur à lacets pur commute à
$\mathfrak{S}_3$. Le dossier ne tranche pas, et nous non plus\footnote{La page
écrit « si $\hat{\hat\Gamma}^{!} \to \boldsymbol{\Gamma}$ est une égalité » ;
la flèche est l'inclusion $\boldsymbol{\Gamma} \subset \hat{\hat\Gamma}^{!}$.
L'exposant $!$ n'est pas défini sur la page ; il est lu d'après la suite
(pages 23 à 25), où $\hat{\hat{\mathcal{E}}}^{!}$ est décrit sans permutation.
L'action du groupe de Galois de $\overline{\mathbb{Q}}/\mathbb{Q}$ tombe dans
$\boldsymbol{\Gamma}$, puisque les $\sigma_i$ sont définis sur $\mathbb{Q}$.}.

\emph{Description symétrique.} Pour décrire $\hat{\hat{\mathcal{E}}}^{!}$ sans
privilégier $L_0$, on donne $u$ par son multiplicateur et trois éléments :
\[
(19)\quad p = \chi(u) \in \widehat{\mathbb{Z}}^{*}, \qquad
(20)\quad g_i \in \widehat{\pi}, \qquad
(21)\quad u(l_i) = \mathrm{int}(g_i)(l_i^{p}) \quad (i = 0, 1, \infty),
\]
liés par
\[
(22)\qquad \mathrm{int}(g_\infty)(l_\infty^{p})\cdot
\mathrm{int}(g_1)(l_1^{p})\cdot \mathrm{int}(g_0)(l_0^{p}) = 1,
\]
et par la condition que les trois facteurs engendrent
$\widehat{\pi}$\footnote{Aux pages 23 et 24 la page écrit
$\mathrm{int}(g_i^{-1})$, l'exposant $-1$ ajouté après coup d'un crayon plus
noir ; aux pages 32 à 37 il a disparu. Les formules (23) et (24), et la phrase
sur $\mathrm{int}(g)\circ u$, ne sont justes que \emph{sans} l'exposant ; on
adopte donc partout la convention des pages 32 et suivantes.}. Pour $u$ donné,
$p$ est unique et chaque $g_i$ est défini à multiplication à droite près par
un élément de $L_i$ ; composer $u$ avec $\mathrm{int}(g)$ garde $p$ et remplace
les $g_i$ par les $g\,g_i$\footnote{La page dit « à multiplication à gauche
près » ; c'est à gauche dans sa convention avec $g_i^{-1}$, à droite dans
celle-ci.}. Si $v$ est donné par $q$ et $(h_i)$,
\[
(23)\qquad vu(l_i) = \mathrm{int}\bigl(v(g_i)\, h_i\bigr)(l_i^{pq}).
\]

\emph{L'extension $S\hat{\hat{\mathcal{E}}}^{!}$.} Soit
$S\hat{\hat{\mathcal{E}}}^{!}$ le groupe des systèmes $(p, g_0, g_1, g_\infty)
\in \widehat{\mathbb{Z}}^{*} \times \widehat{\pi}^{\{0, 1, \infty\}}$ satisfaisant
(22) et la condition de génération, avec la loi que donne (23). C'est une
extension de $\hat{\hat{\mathcal{E}}}^{!}$ par $\widehat{\mathbb{Z}}^{\{0,1,\infty\}}$
(le choix des $g_i$ modulo $L_i$), et les trois fonctions $\gamma_i :
(u, (g_j)) \mapsto g_i$ satisfont
\[
(24)\qquad \gamma_i(\tilde v\tilde u) = \tilde v\bigl(\gamma_i(\tilde u)\bigr)\,
\gamma_i(\tilde v) :
\]
$(\gamma_i)$ est un $1$-cocycle de $S\hat{\hat{\mathcal{E}}}^{!}$ à valeurs dans
$\widehat{\pi}^{\{0,1,\infty\}}$\footnote{Avec le produit dans cet ordre, c'est
$(\gamma_i^{-1})$ qui satisfait la règle usuelle $c(vu) = c(v)\,v(c(u))$ des
homomorphismes croisés pour l'action à gauche. Deux premières écritures du
second membre de (24) sont biffées sur la page.}. Le noyau de
$S\hat{\hat{\mathcal{E}}}^{!} \to \mathrm{Out}^{!}_{\mathrm{lac}}(\widehat{\pi})$
s'identifie au produit
\[
\widehat{\pi} \times \prod_{i} L_i \simeq \widehat{\pi} \times
\widehat{\mathbb{Z}}^{\{0,1,\infty\}}, \qquad (g, (\alpha_i)) \longmapsto
\bigl(1, (g\, l_i^{\alpha_i})_i\bigr).
\]

La page 25 se donne pour tâche d'expliciter l'action des $\sigma_i$, des
$\varepsilon_i$ et de $\rho$ sur ces groupes, et en particulier la
« covariance » des $\gamma_i$ par $\rho$. La formule juste est celle que
donnera (15) à la page 36 :
\[
\gamma_{\rho(i)}\bigl(\tilde\rho\,\tilde u\,\tilde\rho^{-1}\bigr) =
\rho\bigl(\gamma_i(\tilde u)\bigr).
\]
C'est la forme que prend la covariance\footnote{La page écrit
$\gamma_i(\rho\tilde u\rho^{-1}) = \gamma_{i+1}(\tilde u)$ : il manque
l'application de $\rho$, et le décalage d'indice est dans l'autre sens. Elle
est d'ailleurs énoncée comme un programme (« il faudrait expliciter \ldots{} et en
particulier \ldots »), non comme un résultat.}.
Une note marginale annonce aussi le groupe $S\widehat{\mathcal{E}}$, extension
de $\widehat{\mathcal{E}}$ par $\widehat{\mathbb{Z}}^{\{0,1,\infty\}}$, et les
relèvements des $\sigma_i$ et de $\rho$ dans ce groupe.

\pagerange{26}{31}
\section*{Scindages partiels et présentations de $\mathcal{E}$ (pages 26 à 31)}

\pagerange{26}{28}
\subsection*{Les sous-groupes finis, point fixe par point fixe (pages 26 à 28)}

Les classes de conjugaison de scindages partiels maximaux de $\mathcal{E}$ —
de ses sous-groupes finis maximaux — correspondent aux points de $U$ fixés par
un élément non trivial de $\Gamma$, et aux composantes connexes des lieux fixes
des involutions antiholomorphes\footnote{La page énumère $\dot\rho\tau$ parmi
ces dernières ; $\dot\rho\tau$ est d'ordre $6$ et sans point fixe, et la page
dit d'ailleurs « éventuellement ».}. La page les passe en revue.

\begin{itemize}
\item \textbf{Les points $P$, $\overline{P}$}, fixes de $\dot\rho$, de
stabilisateur $\Gamma_P$ d'ordre $6$. Le relèvement correspondant est celui des
pages 21 et 20 ; ceux relatifs à $\overline{P}$ s'en déduisent en conjuguant par
un élément de $\mathcal{E}$ au-dessus d'un élément de $\Gamma$ qui envoie $P$
sur $\overline{P}$ ($\dot\sigma_i$, $\tau$, $\tau\dot\rho$, $\tau\dot\rho^2$).
\item \textbf{Les points $0' = 2$, $1' = -1$, $\infty' = \frac12$}, seuls points
fixes dans $U$ de $\dot\sigma_0, \dot\sigma_1, \dot\sigma_\infty$ (l'autre point
fixe, $0$, $1$ ou $\infty$, est ôté). Leurs stabilisateurs $\Gamma_{i'} =
\langle \dot\sigma_i, \tau\rangle \simeq \mathbb{Z}/2 \times \mathbb{Z}/2$ se
relèvent par $\dot\sigma_i \mapsto \sigma_i$, $\tau \mapsto \tau_i$ ; pour
$0'$, le chemin est le quart de grand cercle qui joint $P$ au milieu $2$ de
l'arc $(1, \infty)$ opposé à $0$\footnote{Sur la page, l'indice de
$\Gamma_{\frac12}$ ressemble à un $2$ ; il est lu d'après le point fixe nommé
juste avant. La phrase qui justifie le relèvement (page 27) est très lacunaire.}.
\item \textbf{Le lieu fixe de $\tau$} : $U^{\tau}$ a trois composantes, les arcs
$(0,1)$, $(1,\infty)$, $(\infty, 0)$ du cercle réel, contenant chacun un seul
des points $\infty', 0', 1'$ ; les relèvements correspondants de $\{1, \tau\}$
sont $\tau_\infty$, $\tau_0$, $\tau_1$, et ne donnent rien de nouveau.
\item \textbf{Le lieu fixe de $\tau\dot\sigma_1$}, $z \mapsto 1/\bar z$ : c'est
$\{|z| = 1\} \setminus \{1\}$, connexe, contenant $P$ et $\overline{P}$ ; il
redonne le relèvement $\tilde\sigma_1$ de $\dot\sigma_1\tau$. De même, par
$\rho$, pour $\tau\dot\sigma_\infty$ et $\tau\dot\sigma_0$.
\end{itemize}

\pagerange{28}{30}
\subsection*{Présentations de $\mathcal{E}$ (pages 28 à 30)}

Un premier essai, par un générateur $\rho'$ d'ordre $6$ avec $\rho'^3 = \tau$,
est barré\footnote{Le générateur est écrit sur une rature ; la transcription le
note $\rho'$, les exposants lus « 16 » et « 13 » étant vraisemblablement
$\rho'^6$ et $\rho'^3$. Le paragraphe qui achève la phrase, en haut de la page
29, est barré lui aussi.}. La page 29 part de la présentation de $\Gamma$ par
$\dot\sigma_0, \dot\sigma_1, \dot\sigma_\infty, \dot\rho, \tau$ et relève en
$\sigma_i, \rho, \tau_i$. Les relations vraies dans $\mathcal{E}$ sont
\[
\begin{gathered}
\sigma_i^2 = \tau_i^2 = 1, \qquad \rho^3 = 1, \qquad
\rho\,\tau_i\,\rho^{-1} = \tau_{\rho(i)}, \qquad
\rho\,\sigma_i\,\rho^{-1} = \sigma_{\rho(i)}, \qquad
\tau_i\,\sigma_i = \sigma_i\,\tau_i, \\
\sigma_0\sigma_1 = \rho\,\tau_\infty\tau_1 = \tau_0\tau_\infty\,\rho, \qquad
\sigma_1\sigma_\infty = \tau_1\tau_0\,\rho, \qquad
\sigma_\infty\sigma_0 = \tau_\infty\tau_1\,\rho,
\end{gathered}
\]
la dernière ligne étant (2) récrite par (17), $\rho\, l_0 = l_1\,\rho$\footnote{La
page écrit $\sigma_0\sigma_1 = \tau_\infty\tau_1\rho = l_0\,\rho$, et de même
pour les deux autres : c'est la première écriture, corrigée, des seconds
membres de (2) à la page 17. Elle est fausse dans $\mathcal{E}$ : $l_0\rho \neq
\rho\, l_0$. Les relations $\rho\sigma_i\rho^{-1} = \sigma_{\rho(i)}$ sont
ajoutées en marge (« Il faut ajouter \ldots »), et une seconde note marginale
retrouve (17).}.

La page 30 réduit les générateurs à $\rho, \sigma_0, \tau_0$, les autres s'en
déduisant par conjugaison par $\rho$. La relation qui reste, tirée de
$\sigma_0\sigma_1 = \rho\,\tau_\infty\tau_1$ avec $\sigma_1 = \rho\sigma_0
\rho^{-1}$, $\tau_\infty = \rho^2\tau_0\rho^{-2}$, $\tau_1 = \rho\tau_0\rho^{-1}$,
est
\[
\sigma_0\,\rho\,\sigma_0\,\rho^{-1} = \tau_0\,\rho^{-1}\,\tau_0\,\rho^{-1},
\qquad \text{soit} \qquad (\sigma_0\,\tau_0\,\rho)^2 = 1 ,
\]
et
\[
\mathcal{E} = \bigl\langle\, \rho, \sigma_0, \tau_0 \bigm| \sigma_0^2 =
\tau_0^2 = (\sigma_0\tau_0)^2 = \rho^3 = (\sigma_0\tau_0\rho)^2 = 1
\,\bigr\rangle ,
\]
quotient du produit libre $(\mathbb{Z}/2 \times \mathbb{Z}/2) * \mathbb{Z}/3$
par une seule relation, comme le dit la page\footnote{La page aboutit à
$[\rho, \sigma_0][\rho^{-1}, \tau_0] = 1$. Cette relation est fausse, et pas
seulement dans $\mathcal{E}$ : son image dans $\Gamma$ est $\dot\rho^{-1} \neq 1$.
Deux glissements y mènent : l'ordre $l_0\rho$ de la page 29, puis, au passage à
« i.e. », un $\rho$ écrit pour $\rho^{-1}$. La relation vraie a été vérifiée
dans le groupe libre ; l'accolade de la page sous $\sigma_0\rho\sigma_0^{-1}
\rho^{-1}$ porte « $\sigma_1$ », là où c'est $\sigma_0\sigma_1$.}. Avec
$a = \tau_0$, $b = \tilde\sigma_0 = \sigma_0\tau_0$ et $c = \tilde\sigma_1 =
\tilde\sigma_0\,\rho$, c'est la présentation du groupe de Coxeter
$\langle a, b, c \mid a^2, b^2, c^2, (ab)^2, (bc)^3 \rangle$ du triangle
d'angles $\pi/2, \pi/3, 0$, c'est-à-dire du groupe modulaire étendu
$PGL_2(\mathbb{Z})$. En effet, les relations valant dans $\mathcal{E}$, on
a un morphisme surjectif de ce groupe de Coxeter sur $\mathcal{E}$ ; il envoie
le sous-groupe d'indice $2$ des mots pairs, engendré par $ab = \sigma_0$ et
$bc = \rho$ et classiquement isomorphe à $\mathbb{Z}/2 * \mathbb{Z}/3$,
isomorphiquement sur $\mathcal{E}^{+}$ (page 19), et il est compatible aux
deux morphismes vers $\mathbb{Z}/2$ ; c'est donc un isomorphisme\footnote{L'identification à $PGL_2(\mathbb{Z})$ est de cette édition.}.

La page remarque enfin une symétrie qui échange $\sigma_0$ et $\tau_0$, et
$\rho$ et $\rho^{-1}$ ; elle respecte bien les relations ci-dessus, puisque
$\tilde\sigma_0$ conjugue $\rho$ en $\rho^{-1}$. Mais elle n'est pas compatible
avec la structure d'extension : si elle préservait $\pi$, elle induirait sur
$\Gamma$ un automorphisme échangeant $\tau$, central, et $\dot\sigma_0$, qui ne
l'est pas\footnote{La page le dit « comme on voit de suite », avec une
parenthèse incertaine sur le passage au quotient ; l'argument par le centre est
de cette édition.}.

\pagerange{31}{31}
\subsection*{L'extension $S\Gamma$ (page 31)}

Il faudrait encore expliciter l'extension $S\Gamma = S\Gamma_{0,3}$ de $\Gamma$
par $\mathbb{Z}^{\{0,1,\infty\}}$, d'où, par produit fibré, $S\mathcal{E} =
S\Gamma \times_\Gamma \mathcal{E}$. La page s'arrête là ; la suite
(pages 32 à 37) construit d'abord le groupe complété et revient à
$S\mathcal{E}$ par restriction.

\pagerange{32}{37}
\section*{Le groupe $S\hat{\hat{\mathcal{E}}}$ (pages 32 à 37)}

\pagerange{32}{33}
\subsection*{Coordonnées et loi de groupe (pages 32 et 33)}

« (2) Le groupe $S\hat{\hat{\mathcal{E}}}$ » : c'est l'ensemble des systèmes
\[
(1)\qquad \tilde u = (p, \sigma, g_0, g_1, g_\infty), \qquad p \in
\widehat{\mathbb{Z}}^{*},\ \sigma \in \mathfrak{S}_3,\ g_i \in \widehat{\pi},
\]
satisfaisant
\[
(2)\qquad \mathrm{int}(g_\infty)(l_{\sigma(\infty)}^{p}) \cdot
\mathrm{int}(g_1)(l_{\sigma(1)}^{p}) \cdot \mathrm{int}(g_0)(l_{\sigma(0)}^{p})
= 1,
\]
et tels que les trois facteurs engendrent $\widehat{\pi}$ ; à $\tilde u$ est
associé l'automorphisme à lacets
\[
(3)\qquad u(l_i) = \mathrm{int}(g_i)(l_{\sigma(i)}^{p}), \qquad
(4)\qquad S\hat{\hat{\mathcal{E}}} \longrightarrow \hat{\hat{\mathcal{E}}} =
\mathrm{Aut}_{\mathrm{lac}}(\widehat{\pi}), \quad \tilde u \mapsto u .
\]
Le morphisme
\[
(5)\qquad i : \widehat{\pi} \times \widehat{\mathbb{Z}}^{\{0,1,\infty\}}
\longrightarrow S\hat{\hat{\mathcal{E}}}, \qquad
(6)\qquad i(g; \alpha_0, \alpha_1, \alpha_\infty) = (1, 1, g\, l_0^{\alpha_0},
g\, l_1^{\alpha_1}, g\, l_\infty^{\alpha_\infty}),
\]
est un homomorphisme injectif du produit direct, et (7) l'image de
$i(g; \alpha)$ par (4) est $\mathrm{int}(g)$. Donc (4) est surjectif, de noyau
$i(\{1\} \times \widehat{\mathbb{Z}}^{\{0,1,\infty\}})$, et
$i(\widehat{\pi} \times \widehat{\mathbb{Z}}^{\{0,1,\infty\}})$ est le noyau du
composé $S\hat{\hat{\mathcal{E}}} \to
\mathrm{Out}_{\mathrm{lac}}(\widehat{\pi})$\footnote{La note marginale dit que
(4) « a comme noyau $i(\hat\pi)$ » ; ce n'est ni l'un ni l'autre de ces deux
noyaux.}. La multiplication est
\[
(8)\qquad (p, \sigma, g_0, g_1, g_\infty)\,(p', \sigma', g'_0, g'_1, g'_\infty)
= \bigl(pp',\ \sigma\sigma',\ (u(g'_i)\, g_{\sigma'(i)})_{i}\bigr),
\]
comme on le vérifie en appliquant $u$ à $u'(l_i) = \mathrm{int}(g'_i)
(l_{\sigma'(i)}^{p'})$. Elle n'est pas « entièrement explicite », note la page,
puisqu'elle fait intervenir $u(g'_i)$, calculable seulement quand $g'_i$ est
donné comme un mot (ou une limite de mots) en les $l_j$. Elle l'est quand le
premier facteur est dans l'image de $i$ :
\[
(9)\qquad i(g; \alpha)\,(p, \sigma, g_i) = \bigl(p, \sigma, (g\, g_i\,
l_{\sigma(i)}^{\alpha_{\sigma(i)}})_i\bigr),
\]
ce qui pour $\sigma = 1$ donne (9') $(p, 1, (g\, g_i\, l_i^{\alpha_i})_i)$ ; et
l'action de $S\hat{\hat{\mathcal{E}}}$ par conjugaison sur l'image de $i$ ne
dépend que de $u$ :
\[
(10)\qquad \tilde u\; i(g; \alpha_0, \alpha_1, \alpha_\infty)\; \tilde u^{-1}
= i\bigl(u(g);\ p\,\alpha_{\sigma^{-1}(0)},\ p\,\alpha_{\sigma^{-1}(1)},\
p\,\alpha_{\sigma^{-1}(\infty)}\bigr),
\]
avec la même réserve sur l'explicitation de $u(g)$\footnote{Les indices de (8)
et (9) sont lus avec peine ; ceux qui sont écrits ici sont ceux que le calcul
impose, et ils coïncident avec la lecture de la transcription. Le point
d'insertion de « mais ils s'explicitent ainsi », page 33, n'est pas sûr.}. En
particulier $i(\widehat{\pi})$ est distingué, et le quotient
$S\hat{\hat{\mathcal{E}}}/i(\widehat{\pi})$ est une extension de
$\mathrm{Out}_{\mathrm{lac}}(\widehat{\pi})$ par
$\widehat{\mathbb{Z}}^{\{0,1,\infty\}}$ dont $S\hat{\hat{\mathcal{E}}}$ est le
produit fibré avec $\hat{\hat{\mathcal{E}}}$ : c'est la forme complétée du
$S\mathcal{E} = S\Gamma \times_\Gamma \mathcal{E}$ de la page 31\footnote{Cette
remarque est de cette édition. Les $g_i$ jouent le rôle de chemins du point
base vers chacune des trois pointes ; la notion de point base tangentiel
(Deligne, 1989), qui les rendrait canoniques, n'est pas dans le dossier.}.

\pagerange{34}{35}
\subsection*{Les rectificateurs triviaux (pages 34 et 35)}

On regarde la partie $S\widehat{\mathcal{E}}$ de $S\hat{\hat{\mathcal{E}}}$ au-dessus
de $\Gamma \subset \hat{\hat\Gamma}$, et les éléments $\rho$, $\sigma_i$, $\tau_i$,
qui sont même dans $\mathcal{E}$. Le relèvement multiplicatif $\dot\rho \mapsto
\rho$, $\dot\sigma_0\tau \mapsto \tilde\sigma_0 = \sigma_0\tau_0$ de $\Gamma_P$
dans $\mathcal{E}$ se remonte à $S\mathcal{E}$ par
\[
(11)\qquad \varphi : \Gamma_P \longrightarrow S\mathcal{E}, \qquad
(12)\qquad
\begin{aligned}
\dot\rho &\longmapsto (1, \rho, 1, 1, 1) =: \tilde\rho,\\
(\dot\sigma_i, \tau) &\longmapsto (-1, \sigma_i, 1, 1, 1) =: \tilde\sigma'_i ,
\end{aligned}
\]
où $\tilde\rho$ est la permutation $l_i \mapsto l_{\rho(i)}$ et
$\tilde\sigma'_i$ l'automorphisme $l_j \mapsto l_{\sigma_i(j)}^{-1}$ — pour
$i = 0$, celui de (14)\footnote{La page note $\Gamma^{*}_P$ ce qui est ici
$\Gamma_P$, et écrit les éléments de $\Gamma_P$ comme $(\dot\sigma, \pm 1)$.}.

L'ensemble des $(p, \sigma, 1, 1, 1)$ est un sous-groupe (par (8)), formé des
$(p, \sigma)$ tels que
\[
(13)\qquad l_{\sigma(\infty)}^{p}\, l_{\sigma(1)}^{p}\, l_{\sigma(0)}^{p} = 1
\]
— c'est (2) avec $g_i = 1$\footnote{La page écrit (13) dans l'ordre inverse,
$l_{\sigma(0)}^{p} l_{\sigma(1)}^{p} l_{\sigma(\infty)}^{p} = 1$ ; pour $\sigma
= 1$, $p = 1$ ce serait $l_0 l_1 l_\infty = 1$, faux sous (1). L'ordre juste est
celui de (2) et de (14).}. Dans $S\mathcal{E}$, où $p = \pm 1$, les solutions
sont exactement les $(\operatorname{sg}\sigma, \sigma)$ : pour $p = 1$, il faut
que $(\sigma(\infty), \sigma(1), \sigma(0))$ soit une permutation circulaire de
$(\infty, 1, 0)$, et pour $p = -1$, que $(\sigma(0), \sigma(1),
\sigma(\infty))$ le soit. Le sous-groupe $\tilde\Gamma_P = \varphi(\Gamma_P)$
est donc caractérisé dans $S\mathcal{E}$ comme celui des automorphismes qui
préservent $L_0, L_1, L_\infty$, affectés des « rectificateurs » $g_0 = g_1 =
g_\infty = 1$.

\emph{L'exercice.} La page conjecture (« Sans doute ») que le même résultat vaut
dans $S\hat{\hat{\mathcal{E}}}$, c'est-à-dire que (13) force $p = \pm 1$, le
ramène — quitte à multiplier par un $\tilde\sigma'_i$ pour rendre $\sigma$
paire, ce qui change $p$ en $-p$ — à l'énoncé : pour $p \in
\widehat{\mathbb{Z}}^{*}$,
\[
(14)\qquad l_\infty^{p}\, l_1^{p}\, l_0^{p} = 1 \quad\Longrightarrow\quad p = 1,
\]
et le laisse en « Exercice ! ». C'est vrai, et voici une solution\footnote{La
solution est de cette édition. La phrase de la page qui opère la réduction
(« Je pense qu'on peut [quitte : corriger par un $\tilde\sigma_i$, ou $\sigma \in
\Gamma_P$)] revient au $\equiv$ de dire\ldots ») est lue avec peine ; la réduction
elle-même est juste, les $\sigma$ paires donnant des conjugués de (14).}. Par
(1), $l_\infty = (l_1 l_0)^{-1}$, et (14) s'écrit $(l_1 l_0)^{p} = l_1^{p}
l_0^{p}$. Si $p \neq 1$, il existe un entier $N \geq 3$ tel que $p \not\equiv 1
\pmod N$. Envoyons $l_0, l_1$ — c'est permis, $\widehat{\pi}$ étant libre profini
sur eux — sur deux réflexions $x, y$ du groupe diédral d'ordre $2N$ dont le
produit $r = yx$ est une rotation d'ordre $N$. Comme $p$ est impair, $x^{p} = x$
et $y^{p} = y$, donc $l_1^{p} l_0^{p} \mapsto r$, tandis que $(l_1 l_0)^{p}
\mapsto r^{p} \neq r$. Donc (14) entraîne $p = 1$, et (13) entraîne $p = \pm 1$
et $(\sigma, p) \in \Gamma_P$ dans $S\hat{\hat{\mathcal{E}}}$ aussi.

\pagerange{36}{37}
\subsection*{Relèvements des $2$-Sylow ; centralisateurs de $\tilde\rho$ et
$\tilde\sigma_0$ (pages 36 et 37)}

Les $2$-sous-groupes de Sylow $\langle \dot\sigma_i, \tau \rangle \simeq
\mathbb{Z}/2 \times \mathbb{Z}/2$ de $\Gamma$, qui se relèvent dans
$\mathcal{E}$ (pages 26 à 28), ne se relèvent \emph{pas} dans $S\mathcal{E}$ :
déjà $\dot\sigma_0$ ne se relève pas en un élément d'ordre $2$. L'argument de la
page est que l'extension locale en $0$ n'est pas scindée : aucun élément
d'ordre $2$ de $\mathcal{E}$ au-dessus de $\dot\sigma_0$ ne centralise $L_0$. En
effet le centralisateur de $L_0$ dans $\mathcal{E}^{+} \simeq
PSL_2(\mathbb{Z})$ est le stabilisateur d'une pointe, infini cyclique engendré
par $\varepsilon_0$, sans torsion. Et si $\tilde u = (1, \sigma_0, g_0, g_1,
g_\infty)$ était d'ordre $2$ dans $S\mathcal{E}$, (8) donnerait $u^2 = 1$ et
$u(g_0)\,g_0 = 1$, et $w = \mathrm{int}(g_0^{-1})\,u$ serait un élément de
$\mathcal{E}$ au-dessus de $\dot\sigma_0$, fixant $l_0$, avec $w^2 =
\mathrm{int}\bigl((u(g_0)g_0)^{-1}\bigr) = 1$ — contradiction\footnote{Ce
raccord entre $S\mathcal{E}$ et l'extension locale est de cette édition ; la
page le donne comme connu (« on sait bien »). Elle écrit $\Gamma^{+}_0$ pour le
sous-groupe $\langle\dot\sigma_0\rangle$ et $\mathbb{Z}_0$ pour le quotient de
$\mathbb{Z}^{\{0,1,\infty\}}$ par $\mathbb{Z}^{\{1,\infty\}}$ ; « de Sylow » est
un ajout interlinéaire.}.

Plutôt que de présenter $S\Gamma$ et $S\mathcal{E}$ par générateurs et
relations (« je ne vais pas m'y attarder »), la page donne l'action des
automorphismes intérieurs de $\tilde\rho$ et $\tilde\sigma'_0$ :
\[
\begin{aligned}
(15)\quad & \mathrm{int}(\tilde\rho)(p, \sigma, g_0, g_1, g_\infty) =
\bigl(p,\ \rho\sigma\rho^{-1},\ \rho(g_\infty),\ \rho(g_0),\ \rho(g_1)\bigr),\\
(16)\quad & \mathrm{int}(\tilde\sigma'_0)(p, \sigma, g_0, g_1, g_\infty) =
\bigl(p,\ \sigma_0\sigma\sigma_0,\ \tilde\sigma_0(g_0),\ \tilde\sigma_0(g_\infty),\
\tilde\sigma_0(g_1)\bigr),
\end{aligned}
\]
où, dans (16), $\tilde\sigma_0$ désigne l'automorphisme (14) de $\widehat{\pi}$
sous-jacent à $\tilde\sigma'_0$\footnote{La page écrit $\sigma_0(g_i)$ dans (16),
et le second terme de (15) est surchargé (lu $\dot\rho\sigma\dot\rho^{-1}$ par
analogie). Les deux formules se déduisent de (8).}. D'où :

\begin{enumerate}
\item[a)] $\tilde u$ commute à $\tilde\rho$ si et seulement si $\sigma$ commute
à $\rho$, soit $\sigma \in \mathfrak{S}_3^{+} = \{1, \rho, \rho^2\}$, et
\[
(17)\qquad g_1 = \rho(g_0), \qquad g_\infty = \rho(g_1), \qquad g_0 =
\rho(g_\infty),
\]
où la troisième égalité suit des deux premières, $\rho$ étant d'ordre $3$. Ces
$\tilde u$ correspondent donc aux $(p, \sigma, g_0) \in
\widehat{\mathbb{Z}}^{*} \times \mathfrak{S}_3^{+} \times \widehat{\pi}$ tels que
\[
(18)\qquad \mathrm{int}\bigl(\rho^2(g_0)\bigr)(l_{\sigma(\infty)}^{p}) \cdot
\mathrm{int}\bigl(\rho(g_0)\bigr)(l_{\sigma(1)}^{p}) \cdot
\mathrm{int}(g_0)(l_{\sigma(0)}^{p}) = 1,
\]
et que les trois facteurs engendrent $\widehat{\pi}$ — l'analogue, dans les
nouvelles coordonnées, de la relation encadrée de la page 7.
\item[b)] $\tilde u$ commute à $\tilde\sigma'_0$ si et seulement si $\sigma \in
\{1, \sigma_0\}$ et $g_0 = \tilde\sigma_0(g_0)$, $g_\infty = \tilde\sigma_0(g_1)$
(donc $g_1 = \tilde\sigma_0(g_\infty)$)\footnote{Une note marginale en oblique,
face à b), dit « variante intéressante pour la suite » ; la suite n'est pas
dans le dossier.}.
\end{enumerate}

Le dossier s'arrête sur une phrase qui annonce le vrai problème, et n'est pas
suivie : « À vrai dire, nous nous intéressons plutôt à la commutation dans
$\hat{\hat\Gamma}$, i.e. la commutation modulo automorphismes intérieurs » —
c'est-à-dire au groupe $\boldsymbol{\Gamma}$ des pages 6 à 11 et 23, qu'il
s'agissait de décrire dans les nouvelles coordonnées.

\end{document}
